\documentclass[10pt,a4paper]{article}

\usepackage[margin=16mm,headheight=15pt]{geometry}
\usepackage{fontspec}
\usepackage{kotex}
\usepackage{xcolor}
\usepackage{array}
\usepackage{booktabs}
\usepackage{longtable}
\usepackage{tabularx}
\usepackage{enumitem}
\usepackage{amsmath,amssymb,bm}
\usepackage[most]{tcolorbox}
\usepackage{hyperref}
\usepackage{fancyhdr}
\usepackage{listings}

\setmainfont{Noto Serif CJK KR}
\setsansfont{Noto Sans CJK KR}
\setmonofont{DejaVu Sans Mono}

\hypersetup{
  colorlinks=true,
  linkcolor=blue!50!black,
  urlcolor=blue!50!black,
  pdftitle={Topology-Distilled Global--Local Wind Dynamics Implementation Checklist},
  pdfauthor={Wind3DGS}
}

\definecolor{TodoGray}{RGB}{105,105,105}
\definecolor{ProgressBlue}{RGB}{35,90,170}
\definecolor{DoneGreen}{RGB}{35,130,80}
\definecolor{BlockRed}{RGB}{170,55,55}
\definecolor{KillRed}{RGB}{125,20,25}
\definecolor{LightBlue}{RGB}{235,243,255}
\definecolor{LightGray}{RGB}{248,248,248}
\definecolor{LightRed}{RGB}{255,239,239}
\definecolor{LightGreen}{RGB}{238,250,242}
\definecolor{LightOrange}{RGB}{255,247,233}

\newcommand{\todo}{\textcolor{TodoGray}{\(\square\)}}
\newcommand{\progress}{\textcolor{ProgressBlue}{[{\raisebox{0.1ex}{\(\sim\)}}]}}
\newcommand{\done}{\textcolor{DoneGreen}{\(\boxtimes\)}}
\newcommand{\blocked}{\textcolor{BlockRed}{[!]}}
\newcommand{\killmark}{\textcolor{KillRed}{[K]}}
\newcommand{\code}[1]{\texttt{#1}}
\newcommand{\artifact}[1]{\textcolor{ProgressBlue}{\nolinkurl{#1}}}
\newcolumntype{L}[1]{>{\raggedright\arraybackslash}p{#1}}
\newcolumntype{Y}{>{\raggedright\arraybackslash}X}

\setlist[itemize]{leftmargin=1.45em,itemsep=0.16em,topsep=0.28em}
\setlist[itemize,2]{label=--,leftmargin=1.8em}
\setlist[enumerate]{leftmargin=1.65em,itemsep=0.2em,topsep=0.3em}
\setlength{\emergencystretch}{2.5em}
\renewcommand{\arraystretch}{1.2}
\sloppy

\lstset{
  basicstyle=\ttfamily\small,
  frame=single,
  breaklines=true,
  columns=fullflexible,
  backgroundcolor=\color{LightGray},
  rulecolor=\color{black!20},
  xleftmargin=0.4em,
  xrightmargin=0.4em,
  aboveskip=0.5em,
  belowskip=0.5em
}

\pagestyle{fancy}
\fancyhf{}
\lhead{Topology-Distilled Global--Local Wind Dynamics}
\rhead{Implementation Checklist}
\cfoot{\thepage}

\title{\vspace{-1.7em}\textbf{Topology-Distilled, Error-Triggered\\Global--Local Wind Dynamics}\\[0.4em]\Large 구현 체크리스트}
\author{Wind3DGS}
\date{Created: 2026-08-13 \quad Status baseline: new method reset}

\begin{document}
\maketitle

\begin{tcolorbox}[
  colback=LightBlue,
  colframe=blue!45!black,
  title=고정 Runtime Pipeline,
  fonttitle=\bfseries
]
\begin{center}
\code{static 3DGS + object package}
\(\rightarrow\)
\code{current surface/aero}
\(\rightarrow\)
\code{Global predictor}
\(\rightarrow\)
\code{benefit gate}
\(\rightarrow\)
\code{active Local base + missing-force channels}
\\[0.25em]
\(\rightarrow\)
\code{overlap assembly}
\(\rightarrow\)
\code{Global complement}
\(\rightarrow\)
\code{single Local physical solve}
\(\rightarrow\)
\code{Global corrector}
\(\rightarrow\)
\code{Gaussian transport/render}
\end{center}

이 문서는 새 연구 방향의 구현 상태를 관리하는 source of truth다.
방법 정의와 수식의 source of truth는
\path{ideas/3dgs_topology_distilled_error_triggered_wind_dynamics_2026-08-13.tex}다.
이전 proxy/SH-residual 체크리스트의 완료 상태는 승계하지 않는다.
\end{tcolorbox}

\begin{tcolorbox}[colback=LightRed,colframe=KillRed,title=가장 중요한 실행 순서 계약,fonttitle=\bfseries]
논문 설명의 Module 번호는 이해를 위한 논리 번호다. 실제 구현은 다음 dependency를 따른다.
\begin{center}
\code{1 $\rightarrow$ 2 $\rightarrow$ 3 $\rightarrow$ 4 $\rightarrow$ 5 $\rightarrow$ 6 $\rightarrow$ 7
$\rightarrow$ 10 $\rightarrow$ 9 $\rightarrow$ 8 $\rightarrow$ 11 $\rightarrow$ 12 $\rightarrow$ 13 $\rightarrow$ 14}
\end{center}
Module 7의 analytic Local base aero, missing aero, missing structural을 channel별로 유지한다.
세 channel을 Module 10에서 먼저 보존적으로 조립하고,
Module 9에서 각각의 Global force span을 제거한 뒤, Module 8에서만 합쳐 단일 Local system을 한 번 적분한다.
Patch별로 적분한 뒤 위치를 평균하는 구현은 금지한다.
\end{tcolorbox}

\tableofcontents
\newpage

\section{문서 사용 규칙}

\subsection{상태 표기}

\begin{center}
\begin{tabular}{@{}ll@{}}
\toprule
상태 & 의미 \\
\midrule
\todo & 아직 시작하지 않음 \\
\progress & 구현 또는 검증 진행 중 \\
\done & 자동 테스트와 완료 기준을 모두 통과함 \\
\blocked & 설계 결정, 외부 자산 또는 환경 문제로 막힘 \\
\killmark & kill gate 발생; 원인 해결 전 후속 milestone 중단 \\
\bottomrule
\end{tabular}
\end{center}

\begin{itemize}
  \item 구현 코드만 존재하는 상태는 \done이 아니다. 재현 명령, 자동 검증, 산출물까지 확인되어야 한다.
  \item Manual inspection이 남아 있으면 최대 \progress로 둔다.
  \item \blocked 항목에는 반드시 \textbf{현재 상태, 원인, 다음 결정}을 기록한다.
  \item \killmark가 발생하면 성능 수치를 완화하여 통과시키지 않고 upstream 설계를 고친다.
\end{itemize}

\subsection{자동 목표 실행 프로토콜}

권장 요청 형식은 다음과 같다.

\begin{lstlisting}
[Wind3DGS | code | TD05] 목표: TD05 완료까지 구현-검증 루프 실행.
중지 조건: 동일 error signature 3회, kill gate, 사용자 판단 필요,
          검증 환경/데이터 부재, 승인 필요한 외부 변경.
\end{lstlisting}

\begin{enumerate}
  \item 각 루프는 \code{구현 -> 검증 -> 수정 -> 재검증} 순서로 진행한다.
  \item 오류는 실행 명령, 실패 check, 핵심 traceback/message를 묶은 error signature로 추적한다.
  \item 같은 signature를 고치는 시도는 최대 3회다. 이후에는 실패 근거와 선택지를 보고한다.
  \item Dataset 선택, 라이선스, 장시간 GPU 학습, system package 설치, 완료 기준 변경,
  외부 서비스, destructive operation은 사용자 판단을 받는다.
  \item 의미 있는 루프가 끝나면 해당 work folder의 \path{sessions/}에
  목표, 변경 파일, 검증 명령, 결과, blocker, 다음 액션을 기록한다.
\end{enumerate}

\subsection{Experiment와 코드 naming convention}

기존 \code{M\#\#} 실험과 충돌하지 않도록 새 방법은 \code{TD\#\#} tag를 사용한다.

\begin{tcolorbox}[colback=LightGray,colframe=black!20,title=권장 경로,fonttitle=\bfseries]
\begin{itemize}
  \item Experiment: \path{experiments/TD##_short_name/}
  \item Reusable code: \path{code/wind3dgs/}
  \item Idea/checklist: \path{ideas/}
  \item Training output/checkpoint: experiment output 또는 명시된 artifact store
\end{itemize}
\end{tcolorbox}

Reusable code는 semantic package로 나눈다.

\begin{lstlisting}
code/wind3dgs/
  contracts/      # schemas, frames, units, force ledger
  teacher/        # teacher conversion and label generation
  topology/       # anchor graph and operator distillation
  aero/           # current-surface law and hyper-reduction
  reduced/        # Global/Local bases and integrators
  local/          # patch proposals, assembly, complement, feedback
  learning/       # topology, missing-force, gate models
  runtime/        # typed 14-module engine
  transport/      # affine anchor-to-Gaussian transport
  evaluation/     # metrics, baselines, profiling
\end{lstlisting}

\section{Scope reset와 현재 상태}

\subsection{이번 체크리스트의 runtime contract}

\begin{longtable}{L{0.23\linewidth}L{0.18\linewidth}L{0.51\linewidth}}
\toprule
항목 & Runtime 허용 & 계약 \\
\midrule
\endfirsthead
\toprule
항목 & Runtime 허용 & 계약 \\
\midrule
\endhead
Canonical static 3DGS & 허용 & Target의 기본 입력이다. \\
Prescribed spatial wind query & 허용 & 유체 상태를 푸는 것이 아니라 control/input으로 조회한다. \\
Sparse oriented anchor graph & 허용 & Non-watertight material scaffold이며 target triangle mesh를 요구하지 않는다. \\
Versioned reduced object package & 허용 & Setup에서 한 번 만들고 frame마다 재구성하지 않는다. \\
Target mesh reconstruction & Mainline 금지 & 별도 end-to-end baseline/fallback에서만 허용하고 결과를 구분한다. \\
Full FEM/MPM/CFD/LBM state & 금지 & Teacher와 accuracy upper bound에서만 사용한다. \\
Frame-wise Gaussian optimization & 금지 & Final anchor state로 affine transport한 뒤 standard renderer를 사용한다. \\
RSH coefficient cache & Mainline 제외 & TD06/E2의 representation ablation으로만 평가한다. \\
\bottomrule
\end{longtable}

\subsection{기존 구현의 재사용 후보}

이전 체크리스트의 \done/\progress 상태를 자동으로 옮기지 않는다.
다음 항목은 새 schema와 E0 contract를 다시 통과하기 전까지 재사용 후보일 뿐이다.

\begin{longtable}{L{0.06\linewidth}L{0.31\linewidth}L{0.55\linewidth}}
\toprule
상태 & 후보 & 새 방법에서 필요한 재검증 \\
\midrule
\endfirsthead
\toprule
상태 & 후보 & 새 방법에서 필요한 재검증 \\
\midrule
\endhead
\progress & Static 3DGS PLY loader/camera loader & 새 \code{CanonicalGaussianAsset} schema, 좌표계, dtype, covariance convention 통과. \\
\progress & Synthetic cloth/leaf fixtures와 debug viewer & Oracle strip/flag fixtures로 재사용하되 물리 정확도 완료로 간주하지 않음. \\
\progress & Position/covariance transport prototype & Final total anchor state에서 단 한 번 수행하는 affine MLS, SPD, normal consistency 재검증. \\
\progress & Procedural wind preview & Wind query fixture와 qualitative smoke test로만 사용; teacher 또는 physical solver가 아님. \\
\blocked & \code{gsplat 1.5.3} on GTX 1080 Ti & Compute Capability 6.1 환경의 기존 JIT incompatibility. 논문 rendering은 compatible backend/GPU가 필요. \\
\bottomrule
\end{longtable}

\subsection{현재 상태 baseline}

\begin{itemize}
  \item[\done] 새 아이디어 스케치, 전용 참고문헌, PDF와 bundle이 존재한다.
  \item[\done] 이전 아이디어 문서와 구현 체크리스트는 backup에 보존했다.
  \item[\done] 새 구현 체크리스트와 milestone dependency를 확정한다.
  \item[\todo] 새 방법의 typed schema와 deterministic contract를 코드로 구현한다.
  \item[\todo] Oracle-topology Global MVP부터 새 completion evidence를 생성한다.
\end{itemize}

\section{고정 입력, 출력, 물리량 owner 계약}

\subsection{Runtime 입력과 출력}

\begin{lstlisting}
RuntimeInput
  canonical_gaussians: G0
  object_package: versioned topology/operators/bases/cubature/binding
  wind_query: W(x, t)
  material_controls
  attachment_controls
  metric_scale_or_nondimensional_preset
  dt

FrameOutput
  deformed_gaussians: Gt
  global_state: q, qdot
  local_state: z, zdot
  patch_state: active_set, decay_set, dormant_set
  diagnostics
  optional_rendered_frame
\end{lstlisting}

\subsection{Force와 state owner ledger}

\begin{longtable}{L{0.22\linewidth}L{0.28\linewidth}L{0.42\linewidth}}
\toprule
물리량 & 유일한 owner & 구현 계약 \\
\midrule
\endfirsthead
\toprule
물리량 & 유일한 owner & 구현 계약 \\
\midrule
\endhead
Mass/inertia & Reduced Global/Local solver & Network가 inertia 또는 acceleration 전체를 다시 예측하지 않는다. \\
Restoring force & \(K_{gg},K_{\ell\ell}\), 실제 nonzero cross block & Gate와 무관하게 항상 존재한다. \\
Damping & \(C_{gg},C_{\ell\ell}\), 실제 nonzero cross block & Gate-off에서도 Local energy를 물리적으로 감쇠한다. \\
Gravity & Explicit external force & Teacher residual에서 정확히 한 번 제거한다. \\
Attachment & Constraint 또는 compliant force 중 하나 & 같은 support를 projection과 force로 이중 적용하지 않는다. \\
Base aerodynamic force & Current-surface analytic law & Relative wind, current normal/area를 사용한다. \\
Missing force & Learned residual network & Base model이 설명한 힘을 제거한 defect만 예측한다. \\
Patch overlap & Constrained force assembly & Network나 displacement averaging으로 seam을 숨기지 않는다. \\
Global complement & Complement projector/basis & Local motion/force의 Global double counting만 제거한다. \\
Global feedback & Corrected-minus-predictor delta & Predictor가 이미 사용한 base/cross term을 다시 넣지 않는다. \\
Gaussian deformation & Final total anchor state & Global/Local이 Gaussian을 각각 두 번 transport하지 않는다. \\
\bottomrule
\end{longtable}

\begin{tcolorbox}[colback=LightOrange,colframe=orange!60!black,title=Gate가 끄는 것과 끄지 않는 것,fonttitle=\bfseries]
Gate는 expensive missing-force expert와 새 Local excitation을 제어한다.
Global/Local mass, stiffness, damping, gravity, attachment, base aero는 끄지 않는다.
Active에서 Decay로 이동할 때 \(z,\dot z\)를 freeze/reset하지 않고 physical decay solve를 유지한다.
\end{tcolorbox}

\section{P0 implementation kill gates}

아래 항목은 단순 checklist가 아니라 후속 milestone 진행을 금지하는 gate다.

\begin{longtable}{L{0.06\linewidth}L{0.45\linewidth}L{0.41\linewidth}}
\toprule
ID & 실패 조건 & 필요한 조치 \\
\midrule
\endfirsthead
\toprule
ID & 실패 조건 & 필요한 조치 \\
\midrule
\endhead
K01 & 좌표계, 단위, normal/front--back, index convention이 artifact마다 다르다. & TD01 schema와 converter를 고치기 전 dataset/package 생성 금지. \\
K02 & \(\|\Phi^\top M\Psi\|\)가 tolerance를 넘는다. & Basis construction/complement를 고치기 전 Local training 금지. \\
K03 & \(M\succ0\), \(K,C\succeq0\), symmetry/cross-block reciprocity가 깨진다. & Runtime solve 금지; operator construction 수정. \\
K04 & Gravity, attachment, aero, restoring, residual에 owner가 둘 이상이다. & Force ledger 수정 후 label과 rollout 재생성. \\
K05 & Gate-off/zero-wind에서 rest 복귀가 실패하거나 Local energy가 증가한다. & Activation, damping, integrator, decay state 수정. \\
K06 & Patch별 integrate-then-average 코드가 있다. & Force assembly 후 단일 Local solve로 구조 변경. \\
K07 & Assembly 후 Global leakage, overlap seam, internal net force/torque가 tolerance를 넘는다. & Joint KKT 또는 basis/patch 수정. \\
K08 & Corrector가 corrected-minus-predictor delta 외 base/cross force를 재사용한다. & Consumed-force ledger와 delta 정의 수정. \\
K09 & Inactive 전환 시 \(z,\dot z\)를 즉시 freeze/reset한다. & Active/Decay/Dormant state와 energy threshold 구현. \\
K10 & Rigid/affine reproduction, covariance SPD, aero/render normal consistency가 실패한다. & Gaussian output/render 금지; transport 수정. \\
K11 & Oracle Local이 mode-count-matched Global-only를 개선하지 못한다. & Local branch 제거 또는 Global/Local basis 재설계. \\
K12 & Always-on missing-force가 physics-only Local을 개선하지 못한다. & Network 제거; learned gate 학습 중단. \\
K13 & Learned gate가 quality--latency Pareto를 개선하지 못한다. & Adaptive claim 제거; always-on 또는 pure physics 채택. \\
\bottomrule
\end{longtable}

\section{Milestone dependency와 전체 현황}

\subsection{Critical path}

\begin{lstlisting}
TD00 governance
  -> TD01 deterministic contracts/E0
  -> TD02 teacher + independent GS data
  -> TD03 oracle object package
  -> TD04 current-surface/full aero
  -> TD05 Global reduced predictor          [MVP-A]
  -> TD06 conservative aero reduction       [MVP-B]
  -> TD07 complementary Local physics       [MVP-C base]
  -> TD08 teacher projection/labels
  -> TD09 always-on missing-force expert
  -> TD10 benefit gate/dispatcher           [MVP-D]
  -> TD11 overlap/complement/feedback
  -> TD12 predicted topology/operator package [MVP-E]
  -> TD13 target runtime + GS transport
  -> TD14 comparison experiments/paper evidence
\end{lstlisting}

TD02의 데이터 생성과 TD04--TD07의 deterministic solver 작업은 interface가 고정된 뒤 병렬 진행할 수 있다.
TD12의 topology network는 Oracle package와 downstream sensitivity가 먼저 확인된 뒤 본격 학습한다.

\subsection{Milestone overview}

\small
\begin{longtable}{L{0.065\linewidth}L{0.235\linewidth}L{0.47\linewidth}L{0.075\linewidth}}
\toprule
ID & Milestone & 핵심 완료 기준 & 상태 \\
\midrule
\endfirsthead
\toprule
ID & Milestone & 핵심 완료 기준 & 상태 \\
\midrule
\endhead
TD00 & Project governance/reset & 새 method용 config, run manifest, artifact 경로와 상태 규칙이 재현된다. & \done \\
TD01 & Deterministic contracts/E0 & Typed schema, force ledger, basis/operator/transport P0 unit tests가 통과한다. & \todo \\
TD02 & Teacher/data pipeline & Force breakdown, split, independent GS가 포함된 누수 없는 dataset이 생성된다. & \todo \\
TD03 & Oracle object package & Oracle anchors/operators/bases/patches/binding package가 validation을 통과한다. & \todo \\
TD04 & Current-surface full aero & Module 1--3 full-sample evaluator가 rigid/relative-wind/normal-area tests를 통과한다. & \todo \\
TD05 & Global ROM predictor & Start/stop/chirp에서 frequency, damping, energy와 same-anchor baseline을 검증한다. & \todo \\
TD06 & Conservative aero reduction & 적은 sample로 generalized force와 net wrench를 함께 보존한다. & \todo \\
TD07 & Complementary Local physics & Oracle active set에서 single Local solve가 stable high-frequency 이득을 보인다. & \todo \\
TD08 & Projection/label generation & Missing-force, gate-benefit, feedback label의 차원과 owner consistency가 확인된다. & \todo \\
TD09 & Missing-force expert & Always-on residual이 physics-only Local보다 rollout/spectrum을 개선한다. & \todo \\
TD10 & Benefit gate/dispatcher & Expert를 실제 skip하며 heuristic보다 같은 budget에서 좋은 Pareto를 만든다. & \todo \\
TD11 & Assembly/complement/feedback & Seam/wrench/double-count/feedback tests를 통과하는 two-way update가 동작한다. & \todo \\
TD12 & Predicted target package & Static GS에서 runtime mesh 없이 topology/operator package를 만들고 Oracle gap을 제한한다. & \todo \\
TD13 & Integrated runtime/transport & 정확한 14-module trace, affine GS transport, p50/p95 profile과 rendering이 재현된다. & \todo \\
TD14 & Evaluation/paper evidence & E0--E10, baselines, OOD, statistics와 go/no-go 결과가 완결된다. & \todo \\
\bottomrule
\end{longtable}
\normalsize

\subsection{Runtime Module과 milestone 대응}

\begin{longtable}{L{0.05\linewidth}L{0.27\linewidth}L{0.21\linewidth}L{0.29\linewidth}}
\toprule
M & 기능 & 주 milestone & 다음 typed 출력 \\
\midrule
\endfirsthead
\toprule
M & 기능 & 주 milestone & 다음 typed 출력 \\
\midrule
\endhead
1 & 이전 state/version 읽기 & TD01, TD13 & \code{PreparedFrameState} \\
2 & Current anchor/surface 복원 & TD04, TD13 & \code{SurfaceState} \\
3 & Relative wind/base aero & TD04 & \code{AeroSampleForces} \\
4 & Conservative aero 축약 & TD06 & \code{ReducedAeroLoad} \\
5 & Always-on Global predictor & TD05 & \code{GlobalPrediction} \\
6 & Future-benefit gate & TD10 & \code{GateDecision} \\
7 & Active Local base+missing-force proposal & TD04, TD09 & \code{PatchForceBatch} \\
10 & Overlap-aware assembly & TD11 & \code{AssemblyResult} \\
9 & Global force span 제거 & TD07, TD11 & \code{ComplementForce} \\
8 & Single Local physical solve & TD07, TD11 & \code{LocalPrediction} \\
11 & Local-to-Global corrector & TD11 & \code{CorrectedDynamicState} \\
12 & Final anchor state & TD13 & \code{FinalAnchorState} \\
13 & Gaussian transport/render & TD13 & \code{DeformedGaussianAsset} \\
14 & State machine/diagnostics & TD10, TD13 & \code{RuntimeState}, \code{FrameDiagnostics} \\
\bottomrule
\end{longtable}

\section{공통 artifact registry}

모든 milestone은 code와 함께 아래 산출물을 남긴다.

\begin{longtable}{L{0.075\linewidth}L{0.31\linewidth}L{0.52\linewidth}}
\toprule
ID & Artifact & 최소 내용 \\
\midrule
\endfirsthead
\toprule
ID & Artifact & 최소 내용 \\
\midrule
\endhead
A00 & \artifact{run_manifest.json} & Git commit, config hash, seed, device, package/dataset version, command. \\
A01 & \artifact{contract_report.json} & Units/frames/shapes, force ledger, PSD/symmetry/orthogonality/transport tests. \\
A02 & \artifact{teacher_manifest.json} & Asset/material/attachment/wind/trajectory/split/force-channel inventory. \\
A03 & \artifact{oracle_object_package/} & Anchors, graph, operators, bases, patches, cubature placeholder, binding, validation. \\
A04 & \artifact{full_aero_report.json} & Relative-wind, normal/area, force component, net wrench unit tests. \\
A05 & \artifact{global_rollout_report.md} & Frequency, amplitude, phase, damping, energy, latency와 failure cases. \\
A06 & \artifact{aero_reduction_report.md} & Full/random/FPS/RSH/direct/proposed accuracy--latency와 wrench. \\
A07 & \artifact{local_physics_report.md} & Basis complement, state transfer, energy, oracle-active Local Pareto. \\
A08 & \artifact{training_labels/} & Projected state, missing-force, gate-benefit, feedback delta, confidence. \\
A09 & \artifact{missing_force_model/} & Config, checkpoint, normalization, validation/rollout report. \\
A10 & \artifact{gate_model/} & Checkpoint, calibration, cost model, oracle regret, dispatcher profile. \\
A11 & \artifact{coupling_report.md} & Overlap seam, force/torque, leakage, corrector double-count와 stability. \\
A12 & \artifact{predicted_object_package/} & Static-GS inferred package, confidence/fallback, Oracle/kNN gap. \\
A13 & \artifact{runtime_report.md} & Exact trace, p50/p95, memory, kernel count, renders, invalid-state counts. \\
A14 & \artifact{paper_evidence/} & E0--E10 tables, plots, videos, statistics, baseline manifests. \\
\bottomrule
\end{longtable}

\newpage
\section{TD00. Project governance와 method reset}

\textbf{목표.} 이전 proxy/SH-residual 방향과 새 Global--Local physics 방향의 상태가 섞이지 않도록
새 experiment namespace, source-of-truth, run manifest와 변경 기록 규칙을 고정한다.

\textbf{선행 조건.} 없음.

\textbf{입력.}
현재 아이디어 스케치, 이 체크리스트, 기존 \path{code/}, \path{experiments/} 자산.

\subsection*{구현 체크}

\begin{itemize}
  \item[\done] 이전 문서와 기존 implementation checklist를 backup에 원래 이름으로 보존한다.
  \item[\done] 현재 아이디어 source/PDF/bibliography/bundle을 \path{ideas/} 최상위에 둔다.
  \item[\done] 새 체크리스트 파일명을 현재 method title에 맞게 고정한다.
  \item[\done] \path{experiments/TD00_contracts/README.md}를 만든다.
  \item[\done] 공통 experiment README template에 다음 필드를 넣는다.
    \begin{itemize}
      \item goal과 falsification question,
      \item exact input privilege와 runtime restrictions,
      \item command/config/seed/device,
      \item outputs/metrics/completion gate,
      \item failed cases와 decision.
    \end{itemize}
  \item[\done] 모든 run이 \artifact{run_manifest.json}을 기록하도록 공통 utility를 만든다.
  \item[\done] config가 dataset/package/model hash를 저장하도록 한다.
  \item[\done] code-side, idea-side, experiment-side session 기록 책임을 README에 명시한다.
  \item[\done] 기존 \code{M01--M03} 구현을 새 완료 상태로 승격하는 audit checklist를 만든다.
  \item[\done] RSH/appearance SH 코드를 main runtime dependency로 import하지 않는지 dependency scan을 추가한다.
\end{itemize}

\subsection*{필수 산출물}

\begin{itemize}
  \item \artifact{experiments/TD00_contracts/README.md}
  \item \artifact{run_manifest.schema.json} 또는 동등한 typed schema
  \item Reuse-candidate audit report
  \item 새 milestone tag를 사용하는 최소 smoke-test run
\end{itemize}

\textbf{완료 기준.}
빈 smoke run 하나가 source commit, config, seed, package/dataset version, device와 output 경로를 자동 기록하고,
이전 method code가 새 completion status에 섞이지 않는다.

\textbf{완료 증거.}
Python 3.12 unit test 38개와 clean-source reference run
\code{td00-contracts-smoke-20260812t184509388327z-abde7ba}가 통과했다.
Compact manifest/report/completion marker의 hash를 독립 verifier로 재검증했다.
Python 3.10/3.11 실제 실행과 물리 E0는 TD01 이후 미검증 항목으로 유지한다.

\section{TD01. Deterministic schema, 좌표계, force ledger와 E0}

\textbf{목표.}
학습을 시작하기 전에 모든 물리량의 shape, dtype, device, 단위, 좌표계, index, owner와 frame dependency를 코드로 고정한다.

\textbf{선행 조건.} TD00 run/artifact convention.

\textbf{입력.}
Canonical GS fixture, 작은 oracle strip fixture, 현재 method의 state/force ownership 정의.

\subsection*{필수 typed data}

\begin{lstlisting}
CanonicalGaussianAsset
TeacherTrajectory
AnchorScaffold
ReducedOperators
AeroCubature
GaussianBinding
ObjectPackage
RuntimeState
PreparedFrameState
SurfaceState
AeroSampleForces
ReducedAeroLoad
GlobalPrediction
GateDecision
PatchForceBatch
AssemblyResult
ComplementForce
LocalPrediction
CorrectedDynamicState
FinalAnchorState
DeformedGaussianAsset
FrameDiagnostics
FrameOutput
\end{lstlisting}

\subsection*{Schema와 convention 체크}

\begin{itemize}
  \item[\todo] World, object, material anchor, local patch, camera frame의 transform 방향을 문서화한다.
  \item[\todo] Right-handed convention, front/back, outward normal과 two-sided sheet sign을 고정한다.
  \item[\todo] Length/time/mass/force 단위를 SI로 둘지 nondimensional preset으로 둘지 artifact마다 명시한다.
  \item[\todo] Metric-scale asset과 normalized asset의 혼용을 validator가 거부하도록 한다.
  \item[\todo] Tensor layout, dtype, device, batch/anchor/patch/sample index convention을 고정한다.
  \item[\todo] Quaternion order와 covariance parameterization을 고정한다.
  \item[\todo] Package, state, model, dataset version과 backward compatibility 규칙을 정의한다.
  \item[\todo] Save--load round trip 후 값, shape, dtype, index, units가 동일한지 검사한다.
  \item[\todo] State와 package version/hash mismatch는 즉시 fail-fast 한다.
  \item[\todo] \code{active\_set}, \code{decay\_set}, \code{dormant\_set}, diagnostics index가 서로 섞이지 않게 한다.
\end{itemize}

\subsection*{Force ledger 체크}

\begin{itemize}
  \item[\todo] 한 frame에 들어오는 모든 force channel의 owner, frame, sample/anchor/generalized space를 기록한다.
  \item[\todo] Gravity, attachment, base aero, structural base, missing force, feedback delta를 enum/channel로 구분한다.
  \item[\todo] Constraint로 소비된 force와 다음 corrector가 소비할 delta를 기록한다.
  \item[\todo] 같은 channel이 predictor와 corrector에서 두 번 더해지면 test가 실패하도록 한다.
  \item[\todo] Teacher residual label 생성 시 제거한 base channel 목록을 label metadata에 저장한다.
\end{itemize}

\subsection*{E0 deterministic unit tests}

\begin{itemize}
  \item[\todo] \(M\) positive definite, \(K,C\) symmetric positive semidefinite test.
  \item[\todo] Reduced block dimension, symmetry와 reciprocity test.
  \item[\todo] \(\Phi^\top M\Phi\approx I\), \(\Phi^\top M\Psi\approx0\) test.
  \item[\todo] Zero force/zero wind에서 rest가 invariant인지 test.
  \item[\todo] Rigid translation/rotation에서 internal elastic energy가 불필요하게 생기지 않는지 test.
  \item[\todo] Constant force와 impulse에 대한 analytic 작은-system response test.
  \item[\todo] Identity/translation/rotation/affine Gaussian reproduction과 covariance SPD test.
  \item[\todo] Net force/torque 계산의 center convention test.
  \item[\todo] CPU/GPU 또는 float32/float64 reference 차이를 기록한다.
  \item[\todo] 고정 seed에서 같은 trace와 artifact hash가 나오는지 test.
\end{itemize}

\subsection*{권장 API}

\begin{lstlisting}[language=Python]
validate_object_package(package) -> ContractReport
validate_runtime_state(state, package, controls, dt) -> PreparedFrameState
validate_force_ledger(ledger) -> ForceOwnershipReport
run_e0_contract_suite(fixtures, tolerances) -> ContractReport
\end{lstlisting}

\textbf{완료 기준.}
\artifact{contract_report.json}에 모든 P0 deterministic check가 통과하고,
schema/version/owner 오류를 의도적으로 주입한 negative test가 정확히 실패한다.
이 milestone 전에는 teacher label이나 learned model을 production dataset으로 만들지 않는다.

\section{TD02. Teacher, independent GS와 dataset split}

\textbf{목표.}
Global/Local solver, missing-force network, gate와 topology distillation을 학습·검증할 수 있는 누수 없는 teacher dataset을 만든다.

\textbf{선행 조건.} TD01 data schema, units/frames와 force ledger.

\textbf{입력.}
Source thin-shell mesh, material/attachment, prescribed wind program, render cameras와 reconstruction config.

\subsection*{Teacher-S structural trajectory}

\begin{itemize}
  \item[\todo] 첫 fixture를 한쪽 고정 cantilever strip/flag로 고정한다.
  \item[\todo] 후속 asset으로 two-point cloth, curtain strip, paper sheet, leaf를 준비한다.
  \item[\todo] High-resolution nonlinear shell/FEM 또는 검증된 dense MPM backend를 선택한다.
  \item[\todo] Timestep, substep, integrator, material law와 damping convention을 manifest에 기록한다.
  \item[\todo] Position, velocity, acceleration을 저장한다.
  \item[\todo] Strain, curvature, stretch/bending energy를 저장한다.
  \item[\todo] Surface normal, area, deformation gradient를 저장한다.
  \item[\todo] Internal, aerodynamic, gravity, attachment/reaction force breakdown을 가능한 한 solver에서 직접 저장한다.
  \item[\todo] 저장된 force breakdown으로 teacher acceleration을 다시 구성하는 balance test를 추가한다.
  \item[\todo] Optional Gaussian mean/covariance target과 multi-view render를 저장한다.
\end{itemize}

\subsection*{Wind/material/attachment sweep}

\begin{itemize}
  \item[\todo] Wind direction/elevation과 weak--strong magnitude sweep.
  \item[\todo] Steady, sinusoidal gust, chirp, abrupt start/stop.
  \item[\todo] Traveling/localized/nonuniform spatial gust.
  \item[\todo] Density, stretch/shear/bending stiffness, damping, anisotropy.
  \item[\todo] One edge, two points, partial/compliant support, stem attachment.
  \item[\todo] Rest, pre-deformed, initial-velocity condition.
  \item[\todo] 각 factor와 random seed가 trajectory manifest에서 검색 가능하도록 한다.
\end{itemize}

\subsection*{Optional Teacher-FSI track}

\begin{itemize}
  \item[\todo] CFD/LBM 또는 계측 subset의 surface sample/force schema를 Teacher-S와 정렬한다.
  \item[\todo] 같은 geometry에서 quasi-steady와 FSI force를 저장한다.
  \item[\todo] \code{unsteady\_defect = FSI/measurement - quasi\_steady}를 명시적으로 만든다.
  \item[\todo] Teacher-FSI가 없는 sample에는 wake/vortex/history label을 생성하지 않는다.
  \item[\todo] Teacher-S-only와 Teacher-FSI sample을 dataset manifest에서 분리한다.
\end{itemize}

\begin{tcolorbox}[colback=LightOrange,colframe=orange!60!black,title=Claim 제한,fonttitle=\bfseries]
Teacher-FSI 또는 실측 force가 없으면 learned force는 basis truncation, nonlinear structure,
operator/topology distillation defect만 보정한다고 기술한다. Wake, vortex shedding 또는 unsteady aerodynamics를 학습했다고 주장하지 않는다.
\end{tcolorbox}

\subsection*{Independent static-GS variants}

\begin{itemize}
  \item[\todo] Source mesh에 Gaussian을 직접 붙인 oracle/debug asset을 별도 표시한다.
  \item[\todo] Rendered views로부터 target static GS를 독립적으로 reconstruct한다.
  \item[\todo] View count, reconstruction seed, Gaussian density variant를 만든다.
  \item[\todo] Hole, floater, covariance noise, opacity perturbation, missing backside variant를 만든다.
  \item[\todo] Teacher와 static GS의 coordinate/metric scale alignment를 저장한다.
  \item[\todo] Alignment residual이 tolerance를 넘는 asset은 fail/fallback으로 표시한다.
\end{itemize}

\subsection*{Split와 leakage 방지}

\begin{itemize}
  \item[\todo] Frame-random split을 금지하고 trajectory 단위 split을 구현한다.
  \item[\todo] Object-disjoint split을 별도로 만든다.
  \item[\todo] Unseen wind/material/attachment/geometry/category/reconstruction split을 만든다.
  \item[\todo] Combined object+material+wind OOD split을 만든다.
  \item[\todo] 같은 source trajectory의 GS variants가 train/test에 흩어져 identity leakage를 만들지 않게 group split한다.
\end{itemize}

\subsection*{필수 산출물}

\begin{lstlisting}
datasets/
  manifests/teacher_manifest.json
  teacher_structure/<trajectory_id>/
  teacher_fsi/<trajectory_id>/          # optional
  canonical_gs/<asset_id>/
  gs_variants/<asset_id>/<variant_id>/
  alignment/<asset_id>/<variant_id>.json
  splits/{train,val,test,ood_*.json}
\end{lstlisting}

\textbf{완료 기준.}
최소 flag/strip의 start/stop, sinusoidal, chirp, traveling gust trajectory가 force breakdown과 함께 재생 가능하고,
force balance, coordinate alignment, trajectory/object split leakage test가 통과한다.

\section{TD03. Oracle runtime object package와 reduced bases}

\textbf{목표.}
Predicted topology의 오차를 배제한 clean oracle case에서 runtime data structure와 Global/Local physical space를 먼저 닫는다.

\textbf{선행 조건.} TD01 contracts, TD02 oracle source/trajectory.

\textbf{입력.}
Source mesh topology와 material/attachment, canonical/independent GS, teacher snapshots.

\subsection*{Oracle AnchorScaffold}

\begin{itemize}
  \item[\todo] Material-space anchor 위치, rest normal, area, mass와 reliability를 만든다.
  \item[\todo] Oriented material edges와 bending/hinge relations를 저장한다.
  \item[\todo] Hard/compliant attachment를 하나의 명시적 representation으로 저장한다.
  \item[\todo] Core--halo patch와 overlap graph를 만든다.
  \item[\todo] Partition-of-unity로 total area와 mass가 원본과 일치하는지 검사한다.
  \item[\todo] Disconnected close layer를 일부러 포함한 negative fixture를 만든다.
\end{itemize}

\subsection*{Physical operators}

\begin{itemize}
  \item[\todo] Positive diagonal/lumped \(M\)을 조립한다.
  \item[\todo] Stretch, shear, bending, attachment \(K\) block을 조립한다.
  \item[\todo] Rayleigh 또는 명시적 local damping \(C\)를 조립한다.
  \item[\todo] PSD, symmetry, attachment null-space와 rigid-mode test를 실행한다.
  \item[\todo] Material/control 변경 시 어떤 block과 factorization을 invalidate하는지 기록한다.
\end{itemize}

\subsection*{Global/Local bases}

\begin{itemize}
  \item[\todo] Generalized eigenmode 또는 physics/POD hybrid로 \(\Phi\)를 만든다.
  \item[\todo] \(\Phi^\top M\Phi=I\)를 검사한다.
  \item[\todo] Patch raw basis \(B_r\)와 patch-to-global scatter를 만든다.
  \item[\todo] Patch basis를 anchor space에 먼저 overlap-aware assemble한다.
  \item[\todo] Assemble된 raw basis에서 Global span을 제거하고 \(\Psi\)를 \(M\)-orthonormalize한다.
  \item[\todo] \(\Phi^\top M\Psi=0\) tolerance를 asset/rank별로 저장한다.
  \item[\todo] \(M_{gg},C_{gg},K_{gg},M_{\ell\ell},C_{\ell\ell},K_{\ell\ell}\)와 cross block을 만든다.
  \item[\todo] Cross block이 수치적으로 0인지 nonzero인지 기록하고 feedback claim과 연결한다.
  \item[\blocked] Dynamic active basis와 fixed-superset sparse mask 중 첫 구현 방식을 결정한다.
    \begin{itemize}
      \item 권장 첫 구현: fixed superset \(\Psi\) + sparse patch excitation mask.
      \item Dynamic basis는 state remapping/energy continuity를 별도 통과한 뒤 사용한다.
    \end{itemize}
\end{itemize}

\subsection*{Gaussian binding}

\begin{itemize}
  \item[\todo] 각 Gaussian을 고정된 소수 anchor support와 partition weight에 bind한다.
  \item[\todo] Rest offset와 optional local frame을 저장한다.
  \item[\todo] Frame마다 nearest anchor를 다시 고르지 않게 한다.
  \item[\todo] Identity/translation/rotation/affine reproduction을 TD01 suite와 연결한다.
\end{itemize}

\subsection*{권장 API와 package}

\begin{lstlisting}[language=Python]
build_oracle_object_package(
    source_surface, canonical_gs, material, attachment, config
) -> ObjectPackage

save_object_package(package, path)
load_object_package(path) -> ObjectPackage
validate_object_package(package) -> ContractReport
\end{lstlisting}

\begin{lstlisting}
oracle_object_package/
  package_manifest.json
  anchors.npz
  topology.npz
  operators.npz
  bases.npz
  patches.npz
  aero_samples_placeholder.npz
  gaussian_binding.npz
  validation.json
\end{lstlisting}

\textbf{완료 기준.}
Oracle strip package가 save--load 후 동일하며, operator PSD/symmetry, basis orthogonality,
area/mass conservation, binding reproduction test를 모두 통과한다.

\section{TD04. Current-surface state와 full aerodynamic reference}

\textbf{목표.}
Runtime Module 1--3을 learning/hyper-reduction 없이 구현하여 current geometry와 relative wind에서 full-sample aerodynamic force를 얻는다.

\textbf{선행 조건.} TD03 oracle package, TD01 transport/force convention.

\subsection*{Module 1: state preparation}

\begin{itemize}
  \item[\todo] \(q,\dot q,z,\dot z\), active/decay state, memory와 package version을 읽는다.
  \item[\todo] Dimension/index/device/version mismatch를 fail-fast 한다.
  \item[\todo] Moving support/controls를 현재 frame convention으로 변환한다.
  \item[\todo] Wind 값 변경만으로 basis, patch, binding을 재생성하지 않는다.
  \item[\todo] Material/attachment 변경 시 필요한 reduced cache만 invalidate한다.
\end{itemize}

\subsection*{Module 2: current surface reconstruction}

\begin{itemize}
  \item[\todo] \(x=x^0+\Phi q+\Psi z\), \(v=\Phi\dot q+\Psi\dot z\)를 한 번 lift한다.
  \item[\todo] Anchor/sample의 rank-2 또는 affine deformation map을 계산한다.
  \item[\todo] Cofactor 기반 current area-normal을 계산한다.
  \item[\todo] Singular value clamp, orientation continuity와 degeneracy flag를 구현한다.
  \item[\todo] Rest/rigid rotation/in-plane scale/shear fixture로 normal과 area를 검증한다.
  \item[\todo] Aero normal과 rendering transport normal이 같은 convention을 쓰는지 검사한다.
\end{itemize}

\subsection*{Module 3: relative wind와 analytic surface force}

\begin{itemize}
  \item[\todo] Spatial wind를 current sample position에서 query한다.
  \item[\todo] Relative wind를 \(W(x,t)-v\)로 정의한다.
  \item[\todo] Two-sided thin-sheet drag, friction, optional lift term을 구현한다.
  \item[\todo] Current normal/area와 material/aero coefficient를 사용한다.
  \item[\todo] Zero-relative-wind에서 force가 정확히 0인지 검사한다.
  \item[\todo] Uniform rigid translation에서 air/object velocity cancellation을 검사한다.
  \item[\todo] Front/back flip과 normal sign에 대한 force continuity를 검사한다.
  \item[\todo] Component별 force와 net force/torque를 debug output으로 남긴다.
\end{itemize}

\subsection*{권장 API}

\begin{lstlisting}[language=Python]
validate_and_prepare_state(state, package, controls, dt)
    -> PreparedFrameState

reconstruct_surface(prepared_state, package)
    -> SurfaceState

evaluate_quasi_steady_aero(surface, wind_query, aero_parameters, time)
    -> AeroSampleForces
\end{lstlisting}

\subsection*{필수 비교 backend}

\begin{itemize}
  \item Full current-surface/relative-wind reference.
  \item Rest normal/area ablation.
  \item Absolute wind instead of relative wind ablation.
  \item Object-average wind instead of spatial query ablation.
\end{itemize}

\textbf{완료 기준.}
Full-sample evaluator가 analytic unit tests와 teacher surface snapshot의 force/net-wrench comparison을 통과하고,
dynamics를 적분하지 않는 isolated aero report가 생성된다.

\section{TD05. Always-on Global reduced physics predictor}

\textbf{목표.}
Oracle topology에서 Local을 끈 채 object-wide low-frequency wind response를 안정적으로 적분한다.

\textbf{선행 조건.} TD03 reduced operators, TD04 full current-surface aero.

\subsection*{구현 체크}

\begin{itemize}
  \item[\todo] Full sample force를 \(\Phi\)로 projection하여 Global generalized load를 만든다.
  \item[\todo] \(M_{gg}\ddot q+C_{gg}\dot q+K_{gg}q=Q_g\)의 reference solver를 구현한다.
  \item[\todo] Implicit Euler, Newmark 또는 선택한 한 integrator를 first reference로 고정한다.
  \item[\todo] Gravity, attachment/handle, base aero와 structural cross term을 force ledger에 기록한다.
  \item[\done] 이전 Local state의 cross term은 predictor에서 소비하고, TD11 corrector는 updated-minus-consumed delta만 소비하도록 계약을 고정했다.
  \item[\todo] Mainline의 optional learned Global generalized-force head는 default off로 둔다.
  \item[\todo] Mode rank 8/16/32/64 sweep와 factorization cache를 구현한다.
  \item[\todo] Material/attachment/control 변경 시 cache invalidation test를 추가한다.
  \item[\todo] Zero wind/zero displacement rest invariance를 검사한다.
  \item[\todo] Wind start/stop에서 자연스러운 acceleration과 damping envelope를 검사한다.
  \item[\todo] Sinusoidal/chirp에서 modal frequency, amplitude, phase를 teacher와 비교한다.
  \item[\todo] Long rollout에서 energy growth/divergence를 검사한다.
\end{itemize}

\subsection*{권장 API}

\begin{lstlisting}[language=Python]
project_full_aero_to_global(aero_forces, surface, operators)
    -> ReducedAeroLoad

solve_global_predictor(
    previous_state, reduced_load, controls, operators, dt
) -> GlobalPrediction
\end{lstlisting}

\subsection*{필수 baseline}

\begin{itemize}
  \item High-resolution teacher.
  \item Same-anchor sparse shell/PD/XPBD.
  \item WPB-style rest-state/uniform-wind modal projection.
  \item Direct MLP/GNN next-state Global.
  \item Recurrent Global state-space network.
  \item Physics Global + optional learned generalized-force residual.
\end{itemize}

\subsection*{필수 지표}

Trajectory/velocity error, tip amplitude/phase, dominant frequency, damping ratio,
modal coordinate error, energy drift/divergence, p50/p95 dynamics latency와 peak memory.

\textbf{완료 기준.}
MVP-A flag/strip에서 frequency와 damping이 teacher/same-anchor baseline과 일치하고,
4--10초 start/stop/chirp rollout이 발산하지 않는다.
Local 구현 전 mode-count Pareto를 저장하여 Global rank 증가만으로 충분한지 확인한다.

\section{TD06. Conservative current-state aerodynamic hyper-reduction}

\textbf{목표.}
Full current-surface force를 매 frame 전부 계산하지 않고도 Global generalized force와 object net force/torque를 함께 보존한다.

\textbf{선행 조건.} TD04 full evaluator, TD05 Global basis/Jacobian과 rollout reference.

\subsection*{Offline snapshot와 cubature}

\begin{itemize}
  \item[\todo] Training deformation, velocity, spatial wind snapshot에서 full sample force를 계산한다.
  \item[\todo] 각 snapshot의 Global generalized force, net force, net torque를 label로 저장한다.
  \item[\todo] Force/torque center를 rest center, current center 중 하나로 고정하고 manifest에 기록한다.
  \item[\todo] Candidate sample support가 orientation, boundary, attachment, high-curvature 영역을 포함하는지 검사한다.
  \item[\todo] Nonnegative empirical cubature sample과 weight를 최적화한다.
  \item[\todo] Generalized force와 net wrench를 joint objective/constraint로 사용한다.
  \item[\todo] Held-out deformation과 traveling gust snapshot으로 overfit을 검사한다.
  \item[\todo] Sample 수별 실제 GPU kernel/gather/project latency를 측정한다.
  \item[\todo] Estimated cubature error/confidence를 runtime diagnostic과 gate feature로 전달한다.
\end{itemize}

\subsection*{Runtime Module 4}

\begin{itemize}
  \item[\todo] Current selected sample의 surface state와 relative wind만 평가한다.
  \item[\todo] Current Global Jacobian/force direction을 사용하여 \(Q_g^{\mathrm{aero}}\)를 만든다.
  \item[\todo] Weighted net force와 net torque를 동시에 계산한다.
  \item[\todo] Full-sample debug backend를 같은 API로 제공한다.
  \item[\todo] Sample identity는 setup 후 고정하며 wind direction마다 anchor를 다시 고르지 않는다.
  \item[\todo] Weight의 음수/폭주, invalid sample, degeneracy를 runtime에서 검사한다.
\end{itemize}

\subsection*{권장 API}

\begin{lstlisting}[language=Python]
fit_aero_cubature(full_snapshots, operators, constraints, config)
    -> AeroCubature

reduce_aero_load(sample_forces, surface, aero_cubature)
    -> ReducedAeroLoad
\end{lstlisting}

\subsection*{E2 비교 backend}

\begin{enumerate}
  \item Full current-surface samples.
  \item Rest-state/uniform-wind WPB-style projection.
  \item Random samples.
  \item FPS samples.
  \item Learned samples without conservation.
  \item Canonical directional RSH/modal cache.
  \item Low-rank state-conditioned RSH.
  \item Direct direction-conditioned MLP/generalized-force backend.
  \item Proposed conservative current-state cubature.
\end{enumerate}

\begin{tcolorbox}[colback=LightOrange,colframe=orange!60!black,title=RSH 결정 규칙,fonttitle=\bfseries]
RSH는 mainline dependency가 아니다. 동일 sample/data/latency budget에서 held-out direction generalization이나
실제 wall-clock 이득을 명확히 제공할 때만 optional backend로 유지한다. 그렇지 않으면 harmonic direction encoding 또는 direct analytic evaluator를 사용한다.
\end{tcolorbox}

\textbf{완료 기준.}
Held-out deformation/spatial gust에서 full evaluator 대비 generalized-force와 net-wrench tolerance를 만족하고,
random/FPS/nonconservative backend보다 quality--latency 또는 drift Pareto가 좋다.
그렇지 않으면 K13과 별도로 hyper-reduction contribution을 supporting detail로 낮춘다.

\section{TD07. Complementary Local physical system과 Oracle active set}

\textbf{목표.}
Global rank가 놓치는 local flutter/fold를 network 없이 먼저 검증하고,
Global complement에서 하나의 Local physical state가 안정적으로 적분되는지 확인한다.

\textbf{선행 조건.} TD03 \(\Psi\)/Local operators, TD05 Global predictor.

\subsection*{첫 구현 범위}

\begin{itemize}
  \item[\todo] 한 개 non-overlap patch와 fixed \(\Psi\)로 가장 작은 Local solver를 만든다.
  \item[\todo] 다음으로 strongly overlapping two-patch fixture를 추가한다.
  \item[\todo] Oracle future complement energy로 active patch를 선택한다.
  \item[\todo] Physics-only Local과 full-resolution Local upper bound를 만든다.
  \item[\todo] Network force는 0으로 두고 analytic/base local load만 사용한다.
\end{itemize}

\subsection*{Complement와 Local base}

\begin{itemize}
  \item[\todo] Force-space complement와 state-space mass orthogonality를 별도 API로 구분한다.
  \item[\todo] Local force proposal을 \(\Psi\) coordinate로 projection하기 전 Global leakage를 측정한다.
  \item[\todo] \(M_{\ell\ell}\ddot z+C_{\ell\ell}\dot z+K_{\ell\ell}z=Q_\ell\)를 구현한다.
  \item[\todo] Global state에서 오는 base aero/external/cross load를 정확히 한 번 넣는다.
  \item[\todo] Gate/activation 값이 \(M,C,K\)를 scale하지 않도록 한다.
  \item[\todo] Activation blend는 새로운 excitation과 optional output cross-fade에만 적용한다.
  \item[\todo] Active/Decay patch를 하나의 Local coordinate system으로 적분한다.
  \item[\todo] Fixed-superset mask에서 inactive coordinate의 base decay 비용을 profile한다.
\end{itemize}

\subsection*{Dynamic basis를 선택할 경우}

\begin{itemize}
  \item[\todo] Old/New basis 사이 mass-weighted minimum-error state transfer를 구현한다.
  \item[\todo] Position, velocity와 kinetic/potential energy jump를 측정한다.
  \item[\todo] 새 coordinate initialization을 zero, boundary-consistent, projected state로 비교한다.
  \item[\todo] State remap이 P0 tolerance를 넘으면 fixed-superset으로 되돌린다.
\end{itemize}

\subsection*{권장 API}

\begin{lstlisting}[language=Python]
project_to_global_complement(assembled_force, operators)
    -> ComplementForce

solve_local_complement(
    previous_local_state, complement_force, global_prediction,
    operators, activation, dt
) -> LocalPrediction
\end{lstlisting}

\subsection*{검증 체크}

\begin{itemize}
  \item[\todo] \(\Phi^\top M\Psi\), force leakage와 modal amplitude inflation.
  \item[\todo] Zero-wind/gate-off energy monotonic decay.
  \item[\todo] Wind stop 후 physical recovery와 damping envelope.
  \item[\todo] Activation/deactivation position, velocity, acceleration, jerk continuity.
  \item[\todo] Global mode-count 증가 대 Oracle Local의 trajectory/spectrum/latency Pareto.
  \item[\todo] Patch-local normal, curvature, strain과 high-frequency PSD.
  \item[\todo] 4--10초 rollout과 timestep \(\times0.5,\times2\) sensitivity.
\end{itemize}

\textbf{완료 기준.}
Oracle Local이 동일 또는 더 작은 latency 증가로 Global rank 확대보다 local spectrum/trajectory를 반복적으로 개선하고,
zero-wind/activation/state transfer test에서 energy injection이나 discontinuity가 없다.
실패하면 K11에 따라 Local branch를 중단하거나 basis를 재설계한다.

\section{TD08. Teacher projection, missing-force, benefit와 feedback label}

\textbf{목표.}
Network가 위치 전체나 base force를 다시 학습하지 않도록 teacher trajectory에서 state와 missing force를 일관된 공간으로 변환한다.

\textbf{선행 조건.} TD02 teacher force channels, TD03 bases/operators, TD05/TD07 deterministic rollout.

\subsection*{Teacher state projection}

\begin{itemize}
  \item[\todo] Teacher surface state를 material anchors로 resample한다.
  \item[\todo] Mass projection으로 \(q^*,\dot q^*,\ddot q^*\)를 계산한다.
  \item[\todo] Global reconstruction을 제거한 complement에서 \(z^*,\dot z^*,\ddot z^*\)를 계산한다.
  \item[\todo] Teacher velocity/acceleration이 있으면 position finite difference보다 우선한다.
  \item[\todo] Position-only 데이터에는 smoothing, weak-form/integral consistency와 rollout supervision을 둔다.
  \item[\todo] Global/Local reconstruction residual과 patch complement energy를 저장한다.
\end{itemize}

\subsection*{Missing-force inverse dynamics}

\begin{itemize}
  \item[\todo] Local inertia, restoring, damping과 cross block을 label 식에 포함한다.
  \item[\todo] Gravity, base analytic aero, explicit attachment를 정확히 한 번 제거한다.
  \item[\todo] Teacher solver가 직접 제공한 force breakdown을 우선 사용한다.
  \item[\todo] Structural, aerodynamic, attachment/reaction channel을 분리한다.
  \item[\todo] Teacher-FSI sample에만 unsteady aero defect를 표시한다.
  \item[\todo] Label noise/confidence와 inverse-dynamics balance residual을 저장한다.
  \item[\todo] Differentiable Local integrator의 multi-step rollout loss와 force label을 교차 검증한다.
\end{itemize}

\begin{tcolorbox}[colback=LightRed,colframe=BlockRed,title=TD09 전에 반드시 해결할 차원 결정,fonttitle=\bfseries]
문서의 inverse dynamics는 Local generalized force \(\Delta Q_\ell^*\)를 정의하지만,
overlap/conservation pipeline은 patch nodal-force proposal \(\Delta f_r^*\)를 요구한다.
다음 중 하나를 명시적으로 고정해야 한다.
\begin{enumerate}
  \item Teacher anchor-force breakdown을 patch force target으로 직접 사용한다.
  \item \(\Psi^\top f=\Delta Q_\ell^*\)와 wrench/complement 조건을 만족하는 constrained minimum-norm lift로 patch target을 만든다.
\end{enumerate}
변환과 null-space regularization 없이 TD09 network를 학습하면 target과 solver 차원이 일치하지 않으므로 \killmark 처리한다.
\end{tcolorbox}

\subsection*{Counterfactual future-benefit label}

\begin{itemize}
  \item[\todo] 같은 initial state/wind에서 Global-only horizon rollout을 만든다.
  \item[\todo] 각 patch 또는 strongly overlapping group에 Oracle Local을 켠 counterfactual rollout을 만든다.
  \item[\todo] Position, velocity, normal, curvature, strain과 spectrum error 감소를 계산한다.
  \item[\todo] Measured 또는 calibrated Local compute cost를 함께 저장한다.
  \item[\todo] Positive error reduction per cost를 continuous benefit target으로 만든다.
  \item[\todo] Short-horizon label과 일부 long-horizon calibration sample을 만든다.
  \item[\todo] Oracle active set, benefit rank, patch group과 horizon을 저장한다.
  \item[\todo] Teacher future state/force는 label에만 쓰고 runtime feature에서는 제거한다.
\end{itemize}

\subsection*{Feedback label}

\begin{itemize}
  \item[\todo] Predictor geometry의 Global aerodynamic load를 저장한다.
  \item[\todo] Oracle/corrected Local geometry의 Global aerodynamic load를 저장한다.
  \item[\todo] 두 load의 차이 \(\Delta Q_g^{\mathrm{aero},*}\)를 저장한다.
  \item[\todo] 실제 nonzero cross block의 updated-vs-consumed structural delta를 저장한다.
  \item[\todo] Same-frame, next-frame, no-feedback reference를 구분한다.
\end{itemize}

\subsection*{필수 산출물}

\begin{lstlisting}
projected_states/
  q_star, qdot_star, qddot_star
  z_star, zdot_star, zddot_star
  projection_error, complement_energy

missing_force_labels/
  patch_force_channels
  local_generalized_force
  base_channels_removed
  confidence

gate_labels/
  benefit, oracle_active, measured_cost
  error_components, horizon, patch_group

feedback_labels/
  predictor_load, corrected_load, aero_delta
  structural_consumed, structural_updated, structural_delta
\end{lstlisting}

\textbf{완료 기준.}
Label에서 base force를 다시 더하면 teacher balance가 깨지는 negative test와,
runtime-observable feature만으로 dataset을 읽는 leakage test가 통과한다.
Patch-force/generalized-force 차원 결정과 feedback consumed-force ledger가 문서와 코드에 고정된다.

\section{TD09. Always-on active-patch missing-force expert}

\textbf{목표.}
Gate를 붙이기 전에 모든 oracle-relevant patch에서 missing-force expert를 실행하여 network 자체가 physics-only Local보다 유용한지 검증한다.

\textbf{선행 조건.} TD07 stable Local physics, TD08 force target 결정/labels.

\subsection*{입력 feature}

\begin{itemize}
  \item[\todo] Global prediction \(q,\dot q\)와 local history \(z,\dot z\).
  \item[\todo] Patch current position/frame, normal, area, strain, curvature와 deformation-rate invariant.
  \item[\todo] Spatial/relative wind와 base aero force.
  \item[\todo] Material/attachment controls와 topology confidence.
  \item[\todo] Basis truncation/cubature error indicator와 optional recurrent memory.
  \item[\todo] 모든 vector/tensor를 canonical 또는 co-rotated patch frame으로 표현한다.
\end{itemize}

\subsection*{출력 계약}

\begin{itemize}
  \item[\todo] Stage 7 orchestration은 Module 3 analytic law로 active/halo Local base aero를 평가한다.
  \item[\todo] \code{PatchForceBatch}에서 analytic base aero, missing aero, missing structural을 별도 channel/owner/record ID로 유지한다.
  \item[\todo] Position, displacement 또는 next state가 아니라 missing force proposal을 출력한다.
  \item[\todo] Structural internal, aerodynamic external, attachment/reaction channel을 분리한다.
  \item[\todo] External resultant와 internal action--reaction semantics를 metadata에 둔다.
  \item[\todo] Confidence/uncertainty와 optional recurrent memory를 출력한다.
  \item[\todo] TD08에서 선택한 patch nodal-force/generalized-force representation을 정확히 따른다.
  \item[\todo] Proposal 단계에서는 적분, overlap 평균, Global projection을 수행하지 않는다.
\end{itemize}

\subsection*{Model과 학습}

\begin{itemize}
  \item[\todo] 가장 작은 MLP/patch graph model을 baseline으로 만든다.
  \item[\todo] Temporal memory 없는 model을 먼저 학습한다.
  \item[\todo] Gust/history subset에서만 GRU/state-space memory의 추가 이득을 평가한다.
  \item[\todo] Force regression, direction/magnitude, channel consistency loss를 구현한다.
  \item[\todo] Differentiable Local solve를 통과한 one-step/rollout state loss를 추가한다.
  \item[\todo] Work/energy injection penalty와 confidence-weighted noisy-label loss를 평가한다.
  \item[\todo] Teacher forcing뿐 아니라 closed-loop predicted-state rollout을 curriculum에 넣는다.
  \item[\todo] Training normalization과 material/wind scale extrapolation을 기록한다.
\end{itemize}

\subsection*{권장 API}

\begin{lstlisting}[language=Python]
predict_patch_missing_forces(
    active_patch_batch, global_prediction, previous_local_state,
    wind_features, material_controls, model
) -> PatchForceBatch
\end{lstlisting}

\subsection*{필수 비교}

\begin{enumerate}
  \item Global-only.
  \item Full-resolution Local physics upper bound.
  \item Reduced Local physics-only.
  \item Reduced Local + learned missing force.
  \item Direct local displacement network.
  \item Learned force without Local physical integration.
  \item PGRD-style all-node/all-patch residual.
\end{enumerate}

\subsection*{검증}

Patch position/normal/curvature, flutter amplitude/frequency, high-frequency PSD,
wind-stop phase/decay, rollout stability, energy growth와 Local expert latency를 측정한다.

\textbf{완료 기준.}
Always-on learned missing force가 동일 Local physics-only보다 여러 seed/trajectory에서 반복적으로 개선하며,
direct displacement보다 wind-stop recovery와 long rollout이 안정적이다.
개선이 없으면 K12에 따라 network와 TD10 learned gate를 제거한다.

\section{TD10. Future-benefit gate, dispatcher와 patch state machine}

\textbf{목표.}
Local expert가 유익한 patch를 runtime-observable feature로 선택하고,
inactive expert evaluation을 실제로 건너뛰면서 activation discontinuity를 막는다.

\textbf{선행 조건.} TD08 gate labels, TD09 useful missing-force expert.

\subsection*{Gate model}

\begin{itemize}
  \item[\todo] Pooled Global state와 cheap patch invariant만 사용하는 lightweight gate를 만든다.
  \item[\todo] Teacher future state, teacher residual/force를 runtime feature에서 제거한다.
  \item[\todo] Continuous cost-normalized benefit, uncertainty와 predicted cost를 출력한다.
  \item[\todo] Binary current-error target을 ablation으로만 제공한다.
  \item[\todo] Calibration loss/temperature calibration과 uncertainty threshold를 구현한다.
  \item[\todo] Gate feature construction, model inference, selection 비용을 별도 profile한다.
\end{itemize}

\subsection*{Budget와 dispatcher}

\begin{itemize}
  \item[\todo] Benefit/cost TopBudget와 maximum active ratio를 구현한다.
  \item[\todo] Strongly overlapping patch를 group 단위로 선택하는 option을 둔다.
  \item[\todo] On/off threshold, hysteresis와 minimum dwell time을 구현한다.
  \item[\todo] Smooth excitation blend \(\gamma\)를 구현하되 Local \(M,C,K\)에는 곱하지 않는다.
  \item[\todo] Inactive patch를 gather/batch에서 제외하여 expert kernel이 실제 launch되지 않게 한다.
  \item[\todo] Small active set에서 gather/scatter와 launch overhead의 break-even을 측정한다.
  \item[\todo] Always-on이 더 빠른 active ratio에서는 runtime dispatcher가 fallback할 수 있게 한다.
\end{itemize}

\subsection*{Active/Decay/Dormant state}

\begin{longtable}{L{0.16\linewidth}L{0.34\linewidth}L{0.42\linewidth}}
\toprule
상태 & 실행 & 전환 계약 \\
\midrule
Dormant & Cheap gate만 실행 & Benefit이 on threshold/budget을 통과하면 Active. Local state는 충분히 작아야 함. \\
Active & Expert + new excitation + Local physical solve & Off threshold와 dwell 조건을 만족하면 Decay. State reset 금지. \\
Decay & Expert off, excitation fade, Local physical decay solve & Energy/state와 predicted benefit이 충분히 작을 때만 Dormant. 재상승하면 Active. \\
\bottomrule
\end{longtable}

\begin{itemize}
  \item[\todo] Active에서 Decay로 갈 때 \(z,\dot z\)를 보존한다.
  \item[\todo] Decay에서 physical energy가 비증가하는지 검사한다.
  \item[\todo] Recurrent memory의 decay/cleanup 정책을 정한다.
  \item[\todo] Toggle count, activation duration, acceleration/jerk를 기록한다.
  \item[\todo] Active/decay/dormant index의 장시간 memory leak를 검사한다.
\end{itemize}

\subsection*{권장 API}

\begin{lstlisting}[language=Python]
select_local_patches(
    global_prediction, patch_features, previous_state,
    budget, gate_model
) -> GateDecision

update_patch_state_machine(
    gate_decision, local_energy, previous_state, thresholds
) -> PatchStateUpdate
\end{lstlisting}

\subsection*{필수 비교와 지표}

Never, Always, matched-ratio Random, curvature, strain, teacher-residual analysis,
Oracle benefit, learned binary error, proposed cost-normalized gate를 비교한다.
High-benefit recall/precision/AUPRC, calibration, oracle regret, active ratio,
실제 expert call 수, kernel 수, p50/p95 latency, toggle/jerk를 보고한다.

\textbf{완료 기준.}
동일 active budget과 동일 p95 latency 두 조건에서 heuristic보다 낮은 dynamics error를 보이고,
Always-on 대비 실제 wall-clock를 줄인다. Classification score만 높고 expert/kernel이 skip되지 않으면 완료가 아니다.

\section{TD11. Conservative overlap, complement와 Local-to-Global corrector}

\textbf{목표.}
여러 active patch의 force를 하나의 물리계로 조립하고,
Global double counting을 제거한 Local update가 corrected geometry를 통해 Global state에 실제로 feedback되게 한다.

\textbf{선행 조건.} TD07 Local physics, TD09 force proposal, TD10 activation/state metadata.

\subsection*{실제 frame 내부 순서}

\begin{enumerate}
  \item Module 7이 active patch별 analytic Local base aero, missing aero, missing structural proposal을 별도 channel로 출력한다.
  \item Module 10이 각 channel을 anchor force space에 scatter하고 channel별 constraint로 overlap을 조립한다.
  \item Module 9가 각 assembled channel에서 Global span을 제거하거나 joint KKT에서 동시에 제한한다.
  \item Module 8이 channel별 Local generalized force를 처음 합쳐 하나의 Local state \(z,\dot z\)를 한 번 적분한다.
  \item Module 11이 corrected aerodynamic delta와 updated-minus-consumed structural cross delta로 Global corrector를 한 번 푼다.
\end{enumerate}

\subsection*{Module 10: overlap-aware force assembly}

\begin{itemize}
  \item[\todo] Patch-to-global scatter와 core--halo partition weight를 구현한다.
  \item[\todo] Analytic base aero, missing aero, missing structural을 assembly 이후까지 서로 다른 owner/lineage로 유지한다.
  \item[\todo] Structural internal channel의 net force/torque가 0이 되도록 constraint를 둔다.
  \item[\todo] External aerodynamic residual은 unique physical resultant를 보존한다.
  \item[\todo] Attachment reaction은 internal/external residual과 별도 channel로 조립한다.
  \item[\todo] Shared anchor action--reaction과 overlap continuity를 constraint로 둔다.
  \item[\todo] Constrained least-squares/KKT reference solver를 만든다.
  \item[\todo] Patch partition/core--halo 폭 변경에 대한 sensitivity를 검사한다.
  \item[\todo] 단순 displacement/force average는 ablation backend로만 제공한다.
\end{itemize}

\subsection*{Module 9: Global force complement}

\begin{itemize}
  \item[\todo] Force-space projector와 displacement/state projector의 mass metric을 명시한다.
  \item[\todo] \(\Phi^\top f_\ell^\perp\approx0\) 또는 대응 generalized leakage tolerance를 검사한다.
  \item[\todo] Global head가 없다면 removed component를 버리되 diagnostic/loss로 기록한다.
  \item[\todo] Corrected-aero delta와 legitimate cross reaction은 complement 제거 대상이 아님을 channel로 보장한다.
  \item[\todo] Wrench와 complement를 순차 projection했을 때 constraint가 다시 깨지는지 검사한다.
  \item[\todo] 충돌할 경우 wrench와 complement를 하나의 joint KKT로 푼다.
\end{itemize}

\subsection*{Module 8: single Local update}

\begin{itemize}
  \item[\todo] Assembly/complement 결과를 channel별 \(Q_{\ell,c}^{\perp}\)로 변환한다.
  \item[\todo] Active와 Decay coordinate를 하나의 reduced system에서 적분한다.
  \item[\todo] Base analytic Local aero와 두 learned defect는 여기서 처음 합치고, cross load는 별도 structural coupling 항으로 소비한다.
  \item[\todo] Patch별 solver/integrator 호출 후 결과를 average하는 code path가 없음을 trace test로 보장한다.
  \item[\todo] Local energy, work, residual과 state transition diagnostics를 저장한다.
\end{itemize}

\subsection*{Module 11: feedback와 corrector}

\begin{itemize}
  \item[\todo] Predictor와 Local-corrected anchor/surface state를 각각 복원한다.
  \item[\todo] Active/halo region 또는 필요 sample에서 corrected aero를 재평가한다.
  \item[\todo] \(\Delta Q_g^{\mathrm{aero}}=Q_g^{\mathrm{corrected}}-Q_g^{\mathrm{predictor}}\)만 corrector에 넣는다.
  \item[\todo] Full base aero가 corrector에 다시 들어가면 실패하는 unit test를 만든다.
  \item[\todo] Nonzero \(K_{g\ell},C_{g\ell}\)일 때만 structural feedback을 사용한다.
  \item[\todo] One-corrector를 default로 구현한다.
  \item[\todo] Under-relaxation, force-delta clamp와 smaller timestep option을 둔다.
  \item[\todo] No/next-frame/same-frame one-corrector/multiple-iteration을 비교한다.
\end{itemize}

\begin{tcolorbox}[colback=LightRed,colframe=BlockRed,title=Structural feedback 소비 계약,fonttitle=\bfseries]
Module 5는 이전 Local state의 cross coupling을 소비하고 그 값을 \code{predictor\_cross\_load}로 기록한다.
Module 11은 updated reaction 전체가 아니라
\code{coupling(updated local state) - predictor\_cross\_load}인
\code{structural\_cross\_delta}만 더한다.
Local state가 predictor에서 변하지 않았다면 delta가 0이어야 한다.
Force ledger와 negative double-count test 없이 TD11을 완료 처리하지 않는다.
\end{tcolorbox}

\subsection*{권장 API}

\begin{lstlisting}[language=Python]
assemble_patch_forces(patch_force_batch, gate_decision, package)
    -> AssemblyResult

project_to_global_complement(assembly_result, operators)
    -> ComplementForce

solve_local_complement(...)
    -> LocalPrediction

correct_global_state(
    global_prediction, local_prediction, predictor_surface,
    wind_query, operators, consumed_force_ledger, dt
) -> CorrectedDynamicState
\end{lstlisting}

\subsection*{필수 검증}

\begin{itemize}
  \item[\todo] Overlap boundary position/velocity/force jump.
  \item[\todo] Internal net force/torque와 object drift/rotation.
  \item[\todo] Global force/state leakage와 modal amplitude inflation.
  \item[\todo] Corrector base-force/cross-force double-count negative test.
  \item[\todo] Local correction 전후 Global modal state와 teacher error 차이.
  \item[\todo] Corrected-aero delta magnitude/phase와 extra latency.
  \item[\todo] Strong follower-force, abrupt gust와 timestep stress test.
\end{itemize}

\textbf{완료 기준.}
Proposed assembly/complement가 simple averaging보다 seam, net wrench와 drift를 줄이고,
one-corrector가 base force를 중복하지 않으면서 적어도 일부 benchmark의 Global response/teacher error를 개선한다.
Feedback 효과가 없으면 two-way contribution을 제거하고 one-way Local corrector로 단순화한다.

\section{TD12. Static-GS topology/operator distillation과 predicted package}

\textbf{목표.}
Training-only source topology privilege를 이용하되, target runtime에는 static 3DGS만으로 sparse oriented material scaffold와 reduced object package를 생성한다.

\textbf{선행 조건.}
TD03 Oracle package와 TD05--TD11에서 downstream-sensitive operator/basis 요구를 확인했을 것.

\subsection*{Surface reliability와 anchor selection}

\begin{itemize}
  \item[\todo] Opacity, covariance surface-likeness, multi-view contribution과 isolation feature를 계산한다.
  \item[\todo] Boundary/curvature/layer evidence와 reconstruction confidence를 추가한다.
  \item[\todo] Raw Gaussian count를 mass/area/importance로 직접 사용하지 않는다.
  \item[\todo] \(K\ll N_G\) material anchor를 coverage와 reliability 기준으로 선택한다.
  \item[\todo] Anchor rest normal/area/mass/material/attachment/reliability를 예측한다.
  \item[\todo] Partition-of-unity로 total area/mass를 보존한다.
\end{itemize}

\subsection*{Oriented topology와 thin-shell relation}

\begin{itemize}
  \item[\todo] Geometry/evidence 기반 candidate edge generator를 만든다.
  \item[\todo] Material adjacency, rest length, tangent, stretch/shear coefficient를 예측한다.
  \item[\todo] Bending/hinge relation과 orientation consistency를 예측한다.
  \item[\todo] 가까운 disconnected layer와 front/back shortcut을 거부한다.
  \item[\todo] Attachment/compliance와 boundary confidence를 예측한다.
  \item[\todo] Core--halo patches와 overlap graph를 생성한다.
  \item[\todo] Edge/hinge/attachment confidence와 fallback reason을 저장한다.
\end{itemize}

\subsection*{Operator/basis/package construction}

\begin{itemize}
  \item[\todo] Arbitrary dense matrix를 network가 직접 출력하지 않고 positive/symmetric construction으로 \(M,C,K\)를 조립한다.
  \item[\todo] Predicted package에서 TD01 PSD/symmetry/null-space test를 다시 실행한다.
  \item[\todo] Global/Local basis와 cross block을 재구성하고 orthogonality를 검사한다.
  \item[\todo] Aerodynamic sample/cubature를 predicted anchors로 transfer 또는 refit한다.
  \item[\todo] Static Gaussian-to-anchor affine binding을 생성한다.
  \item[\todo] Setup package가 target triangle mesh, FEM/MPM particle grid를 참조하지 않는지 trace한다.
  \item[\todo] Package checksum, version, confidence와 runtime compatibility report를 저장한다.
\end{itemize}

\subsection*{권장 API}

\begin{lstlisting}[language=Python]
infer_anchor_scaffold(
    canonical_gs, material_controls, attachment_controls, model, config
) -> AnchorScaffold

build_object_package(
    canonical_gs, anchor_scaffold, controls, setup_models, config
) -> ObjectPackage
\end{lstlisting}

\subsection*{Fallback와 정직한 scope}

\begin{itemize}
  \item[\todo] Low-confidence edge 제거와 conservative stiffness fallback을 구현한다.
  \item[\todo] User hint/category prior를 사용한 결과를 fully automatic 결과와 분리한다.
  \item[\todo] Target mesh extraction fallback은 별도 baseline/flag로 기록한다.
  \item[\todo] Stable area-normal/bending에 watertight triangle mesh가 사실상 필수라면 target-mesh-free claim을 약화한다.
\end{itemize}

\subsection*{E1 비교와 지표}

Oracle topology, proposed predicted topology, Euclidean kNN, FPS+kNN,
surface reconstruction mesh, no orientation, no bending/stretch separation,
no attachment inference를 비교한다.
Edge/hinge F1, layer shortcut, normal consistency, attachment accuracy, area/mass error,
modal frequency/shape와 downstream rollout/aero/stability를 함께 보고한다.

\textbf{완료 기준.}
Independent static-GS variants에서 package validation이 통과하고,
Predicted topology의 downstream dynamics error가 Oracle package 대비 개발 허용 범위 내에 있다.
Topology F1만 높고 dynamics가 실패하면 완료가 아니다.

\section{TD13. Target runtime engine, final Gaussian transport와 rendering}

\textbf{목표.}
Versioned object package와 static GS에서 정확한 14-module dependency로 한 frame을 실행하고,
최종 total deformation을 Gaussian mean/covariance/frame에 한 번 transport한다.

\textbf{선행 조건.}
TD04--TD12의 standalone modules와 contracts.

\subsection*{통합 Runtime API}

\begin{lstlisting}[language=Python]
class WindDynamicsEngine:
    def setup(
        self,
        gs: CanonicalGaussianAsset,
        controls: Controls,
    ) -> ObjectPackage: ...

    def initialize_state(
        self,
        package: ObjectPackage,
    ) -> RuntimeState: ...

    def step(
        self,
        package: ObjectPackage,
        state: RuntimeState,
        wind: WindQuery,
        controls: Controls,
        dt: float,
    ) -> tuple[FrameOutput, RuntimeState]: ...
\end{lstlisting}

\subsection*{Exact typed frame trace}

\begin{lstlisting}
PreparedFrameState
 -> SurfaceState
 -> AeroSampleForces
 -> ReducedAeroLoad
 -> GlobalPrediction
 -> GateDecision
 -> PatchForceBatch
 -> AssemblyResult
 -> ComplementForce
 -> LocalPrediction
 -> CorrectedDynamicState
 -> FinalAnchorState
 -> DeformedGaussianAsset
 -> RuntimeState / FrameDiagnostics / FrameOutput
\end{lstlisting}

\begin{itemize}
  \item[\todo] Debug mode에서 위 trace와 shape/device/version을 frame별 저장한다.
  \item[\todo] Production mode에서 필요한 diagnostic만 비동기/저비용으로 기록한다.
  \item[\todo] Target runtime이 mesh/MPM/FEM/CFD/LBM code path를 호출하지 않는지 trace test한다.
  \item[\todo] Inactive expert가 실제 skip되는지 kernel/call counter로 검사한다.
  \item[\todo] Factorization, sample, binding, model cache lifetime과 invalidation을 검사한다.
  \item[\todo] NaN/Inf, invalid covariance, inverted area, excessive force/energy를 fail-safe diagnostic으로 잡는다.
\end{itemize}

\subsection*{Module 12: final anchor state}

\begin{itemize}
  \item[\todo] Corrected \(q,z\)를 rest anchor에 단 한 번 lift한다.
  \item[\todo] Final position, velocity, deformation gradient, normal, area를 계산한다.
  \item[\todo] Global map 후 Local map을 순차로 Gaussian에 적용하지 않는다.
  \item[\todo] Hard/compliant attachment error와 seam/leakage를 기록한다.
  \item[\todo] Rank/orientation degeneracy를 confidence/fallback에 전달한다.
\end{itemize}

\subsection*{Module 13: affine Gaussian transport와 rendering}

\begin{itemize}
  \item[\todo] 고정 anchor support의 affine MLS로 Gaussian mean을 transport한다.
  \item[\todo] 동일 deformation map으로 covariance \(F\Sigma^0F^\top\)를 계산한다.
  \item[\todo] Covariance SPD/eigenvalue/scale clamp를 구현한다.
  \item[\todo] Co-rotational appearance/view frame을 갱신한다.
  \item[\todo] Identity/translation/rotation/in-plane affine reproduction을 다시 검사한다.
  \item[\todo] Center-only, rigid-only, full-affine backend를 ablation으로 제공한다.
  \item[\todo] Invalid covariance, scale explosion, frame flip과 temporal flicker를 기록한다.
  \item[\todo] Compatible GPU에서 standard Gaussian renderer를 연결한다.
\end{itemize}

\subsection*{Module 14: final state와 diagnostics}

\begin{itemize}
  \item[\todo] Corrected Global/Local state와 Active/Decay/Dormant state를 저장한다.
  \item[\todo] Hysteresis/recurrent memory의 lifetime을 갱신한다.
  \item[\todo] Module별 p50/p95 latency, kernel/call 수와 memory를 기록한다.
  \item[\todo] Force/torque/energy/conservation/corrector residual을 기록한다.
  \item[\todo] Long rollout에서 state/index/memory가 누적되거나 leak되지 않는지 검사한다.
\end{itemize}

\subsection*{Runtime profiling contract}

다음을 개별 timing scope로 보고한다.

\begin{enumerate}
  \item State/surface reconstruction과 wind query.
  \item Aero force와 hyper-reduction.
  \item Global predictor.
  \item Gate feature/model/selection.
  \item Active expert gather/inference/scatter.
  \item Overlap/complement assembly.
  \item Local integrator.
  \item Global corrector.
  \item Gaussian transport.
  \item Rendering.
\end{enumerate}

Warm-up, GPU synchronization, batch 1/multi-object, p50/p95, peak memory와 kernel count를 모두 기록한다.

\textbf{완료 기준.}
Oracle package와 predicted package 양쪽에서 동일 API가 실행되고,
14-module trace, runtime-mesh-free test, affine transport, long rollout, p50/p95 profile과 rendering output이 재현된다.

\section{TD14. 비교 실험, generalization과 paper evidence}

\textbf{목표.}
각 method component가 실제로 필요한지 반증 중심으로 평가하고,
solver 정확도와 target-mesh-free end-to-end 성능을 공정한 track으로 분리한다.

\textbf{선행 조건.} 해당 component의 standalone artifact와 TD13 integrated runtime.

\subsection*{평가 Track}

\begin{longtable}{L{0.14\linewidth}L{0.34\linewidth}L{0.44\linewidth}}
\toprule
Track & 통제 & 검증 목적 \\
\midrule
A: Oracle-topology matched solver & 동일 topology, mass/material/attachment, wind, timestep, 가능한 동일 DOF & Aero/Global/Local/gate/coupling의 solver 차이. \\
B: Static-3DGS end-to-end & 각 방법의 원래 preprocessing/input privilege를 명시 & Topology distillation, runtime mesh 요구, 전체 quality--latency. \\
C: Real captured static GS & Static multi-view/GS와 별도 dynamic reference & 실제 transfer, trajectory/silhouette/frequency, visual plausibility와 failure. \\
\bottomrule
\end{longtable}

\subsection*{E0--E10 실행 체크}

\begin{longtable}{L{0.07\linewidth}L{0.31\linewidth}L{0.52\linewidth}}
\toprule
ID & 실험 & 완료 증거 \\
\midrule
\endfirsthead
\toprule
ID & 실험 & 완료 증거 \\
\midrule
\endhead
E0 & Deterministic contracts & PSD/symmetry/orthogonality/force owner/transport/decay/double-count tests. \\
E1 & Topology/operator distillation & Oracle/proposed/kNN/FPS/mesh baseline의 static+downstream error. \\
E2 & Aero/RSH alternatives & Full/rest/random/FPS/RSH/direct/proposed force, wrench, sample, latency. \\
E3 & Global physics vs neural & Physics/hybrid/direct/recurrent의 rollout, frequency, damping, OOD, latency. \\
E4 & Local physics/missing force & None/full/reduced/hybrid/direct-displacement의 local spectrum, decay, cost. \\
E5 & Gate/adaptive compute & Never/Always/random/heuristic/oracle/learned의 regret, actual skip, Pareto. \\
E6 & Complement/overlap & No complement, post projection, integrate-average, no torque, proposed의 seam/wrench/leakage. \\
E7 & Feedback/corrector & No/structural/aero/next-frame/one/multi corrector의 Global effect, error, latency. \\
E8 & Static-3DGS end-to-end & Matched direct baselines, runtime privilege, dynamics/rendering/latency/memory. \\
E9 & Real static-GS transfer & Marker/feature trajectory, silhouette, tip/frequency/damping, novel-view temporal quality. \\
E10 & Quality--latency/scaling & Rank/sample/local/active/corrector/Gaussian sweep의 Pareto와 break-even. \\
\bottomrule
\end{longtable}

\subsection*{Generalization와 stress test}

\begin{itemize}
  \item[\todo] Unseen yaw/elevation, magnitude와 combined wind OOD.
  \item[\todo] Unseen gust frequency/chirp와 traveling/localized gust.
  \item[\todo] Unseen stiffness/damping/density, attachment와 geometry.
  \item[\todo] Unseen category와 independent GS reconstruction/corruption.
  \item[\todo] Zero-wind relaxation, abrupt start/stop, threshold crossing.
  \item[\todo] Constant strong wind, resonance chirp, traveling gust.
  \item[\todo] 30--60초 rollout과 timestep \(\times0.5,\times2\).
  \item[\todo] Strong OOD에서 graceful failure, covariance/state validity와 fallback.
\end{itemize}

\subsection*{개발 go/no-go}

\begin{itemize}
  \item[\todo] Predicted topology downstream error가 Oracle보다 약 10--15\% 이상 추가 악화되지 않는지 확인한다.
  \item[\todo] Adaptive Local이 Always-on 대비 dynamics error 약 5\% 내외 증가에서 Local 비용 약 2배 절감을 목표로 한다.
  \item[\todo] End-to-end Always-on 대비 약 1.25배 이상의 의미 있는 speedup 또는 same-quality 개선을 확인한다.
  \item[\todo] MVP Global tip frequency 약 5\%, amplitude 약 10\% 이내를 개발 목표로 사용한다.
  \item[\todo] 4--10초 rollout normalized position error가 object scale 약 5\% 내외이며 divergence가 없는지 확인한다.
\end{itemize}

위 수치는 논문 약속이 아니라 개발 중 방향을 중단하거나 단순화하기 위한 임시 기준이다.

\textbf{완료 기준.}
E0--E10 결과가 재현 가능한 manifest와 함께 생성되고,
핵심 C1--C7 claim마다 적어도 하나의 direct evidence, ablation, falsification result와 failure case가 있다.

\section{학습 curriculum과 loss implementation registry}

학습 가능한 component를 처음부터 joint training하지 않는다.
각 stage는 이전 deterministic/Oracle stage의 완료 증거를 요구한다.

\subsection{Stage 순서}

\begin{longtable}{L{0.06\linewidth}L{0.27\linewidth}L{0.49\linewidth}}
\toprule
Stage & 활성 component & 진입/종료 조건 \\
\midrule
\endfirsthead
\toprule
Stage & 활성 component & 진입/종료 조건 \\
\midrule
\endhead
0 & Deterministic contracts/transport & TD01 E0와 P0 K01--K10 중 해당 항목 통과. \\
1 & Oracle topology + Global physics & TD05 Global-only start/stop/chirp/energy 완료. \\
2 & Oracle Local basis/active set + physics-only Local & TD07 Oracle Local이 mode-count-matched Global보다 유효. \\
3 & Always-on missing-force expert & TD08 force target 차원 확정; TD09가 physics-only Local 개선. \\
4 & Benefit gate + adaptive dispatch & K12 통과 후에만 시작; actual expert skip와 Pareto 요구. \\
5 & Predicted topology/operator & Oracle package downstream sensitivity가 확정된 뒤 distill. \\
6 & Independent GS robustness & Reconstruction seed/density/hole/floater/object-disjoint split 평가. \\
7 & End-to-end transport/rendering & 모든 dynamics P0 통과 후 rendering quality/temporal loss를 연결. \\
\bottomrule
\end{longtable}

\subsection{Loss별 구현 체크}

\begin{itemize}
  \item[\todo] \textbf{Topology loss:} material edge/hinge/orientation/attachment/reliability를 분리해 logging한다.
  \item[\todo] \textbf{Area/mass loss:} local 오차뿐 아니라 object total conservation을 포함한다.
  \item[\todo] \textbf{Operator/modal loss:} PSD construction, eigenvalue/mode shape, downstream rollout을 함께 본다.
  \item[\todo] \textbf{Cubature loss:} generalized force, net force, net torque, nonnegative/sparsity와 latency proxy를 분리한다.
  \item[\todo] \textbf{Missing-force loss:} channel-wise force, direction/magnitude, confidence와 inverse-dynamics residual을 기록한다.
  \item[\todo] \textbf{Gate loss:} benefit regression, ranking, calibration, measured cost와 budget regret를 기록한다.
  \item[\todo] \textbf{Rollout loss:} one-step과 multi-step position/velocity/normal/curvature/spectrum을 분리한다.
  \item[\todo] \textbf{Energy/work loss:} unphysical positive work와 wind-stop decay violation을 penalize한다.
  \item[\todo] \textbf{Overlap/complement loss:} seam, Global leakage, internal net wrench를 기록한다.
  \item[\todo] \textbf{Feedback loss:} predictor/corrected aero delta와 Global state error를 기록한다.
  \item[\todo] \textbf{Rendering loss:} dynamics가 통과한 뒤 PSNR/SSIM/LPIPS/temporal consistency를 추가한다.
  \item[\todo] Total loss weight와 schedule을 config/manifest에 저장하고 stage별 active term을 검증한다.
\end{itemize}

\subsection{Training 안정성과 재현성}

\begin{itemize}
  \item[\todo] Train/validation normalization statistic과 OOD clipping 정책을 저장한다.
  \item[\todo] Teacher-forced one-step와 closed-loop rollout batch 비율을 기록한다.
  \item[\todo] Gradient norm, NaN/Inf, energy violation과 rollout divergence를 monitor한다.
  \item[\todo] Learned component마다 최소 3개 seed를 실행할 수 있게 한다.
  \item[\todo] Best checkpoint selection metric을 paper test metric과 분리한다.
  \item[\todo] Test/OOD asset을 checkpoint selection에 사용하지 않는다.
\end{itemize}

\section{기존 방법 baseline 구현 registry}

Official code를 사용할 수 없거나 input contract를 맞추기 위해 재구현하면 논문명을 그대로 method label로 쓰지 않고
\code{-style} 또는 \code{matched reimplementation}이라고 표시한다.

\subsection{직접 또는 matched 정량 비교}

\small
\begin{longtable}{L{0.20\linewidth}L{0.18\linewidth}L{0.24\linewidth}L{0.18\linewidth}}
\toprule
방법 & 유사점 & 핵심 차이 & 필요한 비교 \\
\midrule
\endfirsthead
\toprule
방법 & 유사점 & 핵심 차이 & 필요한 비교 \\
\midrule
\endhead
High-resolution shell/MPM teacher & Wind-driven deformable surface & Runtime 상한이며 실시간 target 아님 & Accuracy upper bound, label/force balance. \\
Same-anchor sparse PD/XPBD & 동일 graph/DOF의 저렴한 physics & Reduced Global/learned sparse Local 없음 & 속도 이득이 단순 낮은 DOF 때문인지 검증. \\
Wind Projection Basis-style & Directional wind를 modal load로 projection & Rest-state/uniform approximation, GS/Local 없음 & Current-state spatial aero와 Global dynamics E2/E3. \\
Gaussian Swaying-style + matched MPM & Current Gaussian normal/area/relative wind aero & Full MPM, Global ROM/gate 없음 & 같은 aero front-end의 accuracy/latency E8. \\
FreeForm-style mesh-free ROM & Point/GS 기반 reduced mechanics & Wind-specific topology/aero/Local gate 없음 & Oracle/predicted Global package와 ROM 비교. \\
MeshGraphNets-style rollout & Flag-like direct learned dynamics & Mesh/all-node neural, physics base 없음 & Physics 대 direct neural OOD/long stability E3. \\
Recurrent Global network & State/wind에서 next Global state & Restoring/damping을 data로 학습 & Memory가 있어도 physics Global보다 유리한지 E3. \\
Hybrid Global residual & Physics + learned generalized force & Global defect까지 학습 & Global network가 정말 필요한지 E3. \\
PGRD-style all-node residual & Physics + learned per-particle residual & All-node/all-time; spatial benefit gate/complement 없음 & Always-on residual 대 adaptive Local E4/E5. \\
Proposed Global-only & 동일 package/Global backbone & Local off & Local 순수 기여 E4. \\
Proposed Always-on Local & 동일 Local expert & Gate off, 비용 상한 & Gate 품질 상한/비용 상한 E5. \\
Proposed Oracle gate & 동일 expert, perfect selection & Runtime 불가능 & Patch selection upper bound E5. \\
Proposed Learned gate & 최종 방법 & --- & Main quality--latency/Pareto E8/E10. \\
\bottomrule
\end{longtable}
\normalsize

\subsection{조건부 subset 비교}

\begin{longtable}{L{0.22\linewidth}L{0.29\linewidth}L{0.41\linewidth}}
\toprule
방법 & 공통화 가능한 subset & 비교 목적 \\
\midrule
DynamicTree & Tree/branch/leaf asset, external-force compact response & Static GS global modal-like response와 runtime mesh/Local 차이. \\
DiffWind & Prescribed spatial wind forward-simulation subset & Grid MPM/LBM 대 reduced one-way wind dynamics. \\
PhysSkin & 동일 external load/reduced deformation subset & Static GS/mesh-free reduced basis와 Gaussian lifting. \\
PhysGaussian & 동일 material/external force subset & MPM 기반 Gaussian dynamics 비용/품질. \\
GausSim & 동일 elastic object/input이 가능한 subset & Fixed hierarchy 대 error-triggered complementary Local. \\
NGFF & 동일 object-level/global-local force setting subset & Rigid-global/local stress 대 deformable Global modes와 sparse gate. \\
\bottomrule
\end{longtable}

공통 input/action/metric을 만들 수 없으면 정량 main table에 억지로 넣지 않고 capability/failure를 정성적으로만 비교한다.

\subsection{인용/설계 경계 중심}

\begin{itemize}
  \item Subspace Condensation: selective local full-space activation과 two-way coupling의 고전 경계.
  \item Trading Spaces: error/oracle 기반 adaptive enrichment의 경계.
  \item Learning Contact Corrections: learned local nonlinear correction의 경계.
  \item Fast Complementary Dynamics: Global/primary space의 complement constraint 경계.
  \item OmniAD: object-level directional aerodynamic SH representation 경계.
  \item Subspace/cloth two-level solvers: global low-frequency/local high-frequency 분할의 경계.
\end{itemize}

\section{Metric, profiling과 통계 계약}

\subsection{Metric registry}

\small
\begin{longtable}{L{0.17\linewidth}L{0.29\linewidth}L{0.46\linewidth}}
\toprule
범주 & Metric & 구현/보고 계약 \\
\midrule
\endfirsthead
\toprule
범주 & Metric & 구현/보고 계약 \\
\midrule
\endhead
Geometry & Object-scale normalized position/velocity RMSE & Frame와 sequence 평균을 구분하고 object scale 정의를 기록. \\
Surface & Normal angular, strain, curvature error & Invalid/degenerate sample 비율과 함께 보고. \\
Dynamics & Tip amplitude/phase, frequency, damping & Start/stop/sinusoidal/chirp별 time response와 PSD 저장. \\
Aero & Sample/generalized-force error & Full evaluator 기준 magnitude/direction/phase 분리. \\
Conservation & Net force/torque, COM/rotation drift & Internal/external channel과 torque center를 명시. \\
Energy & Work--energy residual, max growth, divergence & Zero wind, wind stop, activation event를 별도 보고. \\
Complement & \(\Phi^\top M\Psi\), force leakage & Basis build와 runtime assembly 둘 다 검사. \\
Patch & Boundary position/velocity/force jump & Core--halo와 activation event에 대한 최대값/분포. \\
Gate & AUPRC, calibration, regret, active ratio & 동일 budget/latency 비교와 actual expert skip 포함. \\
Runtime & Module p50/p95, kernel/call, peak memory & Warm-up, GPU sync, batch/device/version 고정. \\
Topology & Edge/hinge F1, shortcut, attachment, area/mass & Static metric와 downstream dynamics를 함께 보고. \\
Gaussian & Invalid covariance, scale explosion, frame flip & Count와 affected Gaussian/sequence 비율. \\
Rendering & PSNR/SSIM/LPIPS, temporal warp/flicker & Physics main table과 분리하되 end-to-end 결과에 포함. \\
Real data & Marker/feature trajectory, silhouette IoU & Ground-truth 한계를 명시하고 synthetic physics와 분리. \\
\bottomrule
\end{longtable}
\normalsize

\subsection{통계와 reporting}

\begin{itemize}
  \item[\todo] Frame을 독립 표본으로 취급하지 않고 먼저 sequence별 statistic을 계산한다.
  \item[\todo] Asset 단위 평균과 asset bootstrap 95\% confidence interval을 보고한다.
  \item[\todo] Learned baseline과 proposed model은 가능한 최소 3개 seed를 보고한다.
  \item[\todo] Mean, median, worst decile와 divergence/failure sequence 비율을 함께 보고한다.
  \item[\todo] p50/p95 latency는 동일 hardware/backend, warm-up, synchronization으로 측정한다.
  \item[\todo] Same-quality latency와 same-latency quality 두 Pareto를 모두 만든다.
  \item[\todo] Preprocessing/setup/training/runtime cost를 분리한다.
  \item[\todo] Input privilege, runtime mesh/grid, external model/data를 main comparison table에 명시한다.
\end{itemize}

\section{권장 artifact와 run directory layout}

\begin{lstlisting}
assets/
  canonical_gs/
  teacher_sources/

datasets/
  manifests/
  teacher_structure/
  teacher_fsi/
  gs_variants/
  projected_states/
  missing_force_labels/
  gate_labels/
  feedback_labels/

packages/
  <asset_id>/<package_version>/
    package_manifest.json
    anchors.npz
    topology.npz
    operators.npz
    bases.npz
    patches.npz
    aero_cubature.npz
    gaussian_binding.npz
    validation.json

models/
  topology/
  missing_force/
  gate/

runs/
  <run_id>/
    config.yaml
    run_manifest.json
    dataset_hash.txt
    package_hash.txt
    checkpoints/
    diagnostics/
    profiler/
    renders/
    report.md
\end{lstlisting}

\begin{itemize}
  \item[\done] 실제 repository split에 맞는 physical path와 artifact storage를 TD00의 \path{manifests/artifact_policy.json}에서 확정했다.
  \item[\done] Large dataset/checkpoint는 Git에 직접 넣지 않고 manifest와 retrieval rule만 version control하도록 TD00 policy를 고정했다.
  \item[\todo] Artifact schema version migration과 old cache invalidation test를 추가한다.
\end{itemize}

\section{Open Decisions와 결정 시한}

\begin{longtable}{L{0.06\linewidth}L{0.32\linewidth}L{0.23\linewidth}L{0.17\linewidth}}
\toprule
상태 & 결정 질문 & 권장 첫 선택 & 결정 시한 \\
\midrule
\endfirsthead
\toprule
상태 & 결정 질문 & 권장 첫 선택 & 결정 시한 \\
\midrule
\endhead
\blocked & Teacher-S solver/library와 material law는 무엇인가? & 검증 가능한 nonlinear shell/FEM; dense MPM은 matched baseline & TD02 전 \\
\blocked & Metric scale를 어떻게 제공하는가? & Known dimension/scale marker; 없으면 명시적 nondimensional preset & TD01 전 \\
\blocked & Reference timestep/integrator는 무엇인가? & 한 implicit integrator를 고정하고 timestep sweep & TD04/TD05 전 \\
\blocked & Local force network target은 nodal patch force인가 generalized force인가? & Teacher force가 있으면 nodal; 없으면 constrained lift & TD08 전 \\
\blocked & Fixed-superset \(\Psi\)인가 dynamic active basis인가? & Fixed superset + sparse excitation mask & TD07 전 \\
\done & Global predictor가 cross coupling을 소비하는가? & Predictor가 소비하고 corrector는 updated-minus-consumed delta만 소비 & TD00 contract에서 확정 \\
\blocked & Missing-force expert backbone은 무엇인가? & Small co-rotated patch MLP/graph model, memory off & TD09 전 \\
\blocked & Gate backbone/cost model은 무엇인가? & Cheap pooled MLP + measured latency calibration & TD10 전 \\
\blocked & Topology distillation backbone은 무엇인가? & Candidate-edge classifier + positive operator construction & TD12 전 \\
\blocked & Aerodynamic cubature solver/constraint tolerance는 무엇인가? & Nonnegative sparse fit + generalized/wrench joint validation & TD06 전 \\
\blocked & Teacher-FSI를 첫 논문에 넣을 것인가? & Optional small subset; Teacher-S claim과 분리 & TD02 dataset freeze 전 \\
\blocked & 논문용 Gaussian renderer/GPU는 무엇인가? & Compatible CUDA backend, CC 6.1 blocker와 분리 & TD13 전 \\
\bottomrule
\end{longtable}

\section{즉시 실행할 첫 작업 순서}

아래 순서는 learned component를 붙이기 전에 연구 방향을 가장 싸게 반증하는 경로다.

\begin{enumerate}
  \item \textbf{TD00--TD01:} schema, force ledger, E0 fixtures와 test command를 만든다.
  \item \textbf{TD02:} 한쪽 고정 strip의 Teacher-S start/stop/chirp와 force breakdown을 생성한다.
  \item \textbf{TD03:} Oracle anchor/operator/\(\Phi,\Psi\)/binding package를 만든다.
  \item \textbf{TD04--TD05:} Full current-surface aero + Global-only ROM을 same-anchor solver와 비교한다.
  \item \textbf{TD06:} Full/random/FPS/RSH/direct/proposed aero representation probe를 끝낸다.
  \item \textbf{TD07:} Oracle active Local physics가 Global mode 증가보다 나은지 먼저 본다.
  \item \textbf{TD11 deterministic subset:} Oracle force로 assembly/complement/feedback contract를 검증한다.
  \item \textbf{TD08--TD09:} Missing-force target을 확정하고 always-on expert의 필요성을 본다.
  \item \textbf{TD10:} Expert가 유용할 때만 benefit gate와 real skip을 구현한다.
  \item \textbf{TD12--TD13:} 마지막에 predicted topology package와 independent GS runtime을 연결한다.
  \item \textbf{TD14:} 각 단계 artifact를 E0--E10 표/그림으로 승격한다.
\end{enumerate}

\begin{tcolorbox}[colback=LightGreen,colframe=DoneGreen,title=첫 논문 최소 구현선,fonttitle=\bfseries]
Attached single-layer flag/cloth strip/paper/leaf, prescribed spatial wind, Teacher-S 중심,
target-runtime sparse oriented anchor graph, Global physics, Local physics + useful missing force,
benefit gate가 만드는 실제 sparse execution, conservative overlap/complement,
corrected-aero feedback, affine Gaussian transport와 standard rendering까지다.
Full tree, self-contact, tearing, free flight, appearance residual과 arbitrary cross-family material generalization은 후속 범위다.
\end{tcolorbox}

\section{완료 정의와 최종 handoff}

각 milestone을 \done으로 바꾸려면 다음 다섯 증거가 모두 있어야 한다.

\begin{enumerate}
  \item Reusable implementation과 versioned interface.
  \item Reproduction command/config/seed/device가 있는 experiment README.
  \item Positive test와 의도적으로 실패하는 negative test.
  \item 정량 report와 최소 한 개의 visual/trace diagnostic.
  \item P0/go-no-go 결과, 실패 사례와 다음 decision을 기록한 session note.
\end{enumerate}

다음은 완료로 보지 않는다.

\begin{itemize}
  \item Network loss만 감소하고 closed-loop rollout이 검증되지 않은 경우.
  \item Gate score만 높고 expert kernel 또는 wall-clock가 줄지 않은 경우.
  \item Topology F1만 높고 downstream dynamics가 실패한 경우.
  \item Rendering이 그럴듯하지만 force, frequency, damping, energy가 맞지 않는 경우.
  \item Patch 결과를 시각적으로 평균해 seam을 숨긴 경우.
  \item Teacher privilege나 target mesh 사용을 manifest에서 숨긴 경우.
  \item Baseline을 original task/input와 다르게 구현하고 차이를 보고하지 않은 경우.
\end{itemize}

\section*{Notes}

\begin{itemize}
  \item 새 체크리스트는 이전 implementation checklist의 formatting과 검증 원칙만 계승한다.
  \item 기존 완료 상태, mesh-proxy extraction 중심 pipeline, mandatory SH residual은 계승하지 않는다.
  \item Global과 Local 모두 최종 motion은 physical integrator가 계산한다. Network는 missing force와 benefit만 근사한다.
  \item Gate는 learned force evaluation을 선택할 뿐 structural base를 끄지 않는다.
  \item RSH는 공력 방향 표현 ablation이며 paper headline이나 main runtime contract가 아니다.
  \item 가장 먼저 반증할 질문은 \emph{Global rank 확대보다 Local이 필요한가}와
  \emph{physics-only Local보다 missing-force network가 필요한가}다.
\end{itemize}

\end{document}
